\section{Theorem 3.1 for smooth projective varieties}\label{sec:smooth-general}

This chapter proves \cite[Theorem~3.1]{FHX} (\texttt{thm:smooth-EP}) for a smooth projective variety over an
algebraically closed field of characteristic $p>0$, in the model of \S\ref{sec:projective-general}: $X=V(I)\subseteq\PP^M$
projectively normal, with smooth affine charts, and the classes $\theta\,c_1(\mathcal O_X(1))$, $\theta\in\R$
(Theorem~\ref{thm:sg-ep-charp}). The proof is the test-ideal argument of \cite[\S3]{FHX}, organized as the proof for
$\PP^n$ (Remark~\ref{rmk:bj-pn-ep-assembly}).
\begin{itemize}
\item The reduction to integral classes and PL obstacles, the envelope criterion and the description of envelopes by
Fubini--Study functions are proved once for an abstract Fubini--Study system (\S\ref{subsec:sg-envelopes}).
\item Divisorial points of $X$ are defined by discrete valuation rings containing the coordinate ring
$T=S[t]$ of the affine cone over $X\times\A^1$, and the divisorial comparison theorem \cite[Theorem~4.22]{BJ22} is
proved on $X$ from Rees valuations on the charts $A_j[t]$ (\S\ref{subsec:sg-comparison}). The ideals $J_m$ are
defined from the obstacle, as on $\PP^n$ (\S\ref{subsec:sg-Jm}).
\item Test ideals are developed for rings with a $p$-basis (\S\ref{subsec:sg-pbasis}); smooth algebras are covered by
principal open sets with $p$-bases by Kunz's argument (\S\ref{subsec:sg-smooth}).
\item On the punctured cone over $X\times\A^1$ the asymptotic test ideals are written globally as
$\Phi^e(J_r^NM_0)\,\mathfrak J_0$, where $\Phi^e$ is the Frobenius trace for a homogeneous $p$-basis of $\operatorname{Frac}T$
and $M_0$, $\mathfrak J_0$ correct for its ramification (\S\ref{subsec:sg-cone}). This form gives the uniform twist of
\cite[Lemma~3.2]{FHX} by a count of degrees (Proposition~\ref{prop:sg-uniform-generation}) and the valuative bound of
\cite[Lemma~3.3]{FHX} from the Frobenius trace estimate (Proposition~\ref{prop:sg-valuative-bound}).
\end{itemize}
Castelnuovo--Mumford regularity, Serre vanishing, Frobenius pushforwards of sheaves, canonical bundles, test
configurations, resolution of singularities, countability of $k$ and \cite[Theorem~6.21]{BJ22} are not used.
Remark~\ref{rmk:sg-source} compares the route with \cite[\S3]{FHX}. In characteristic $0$ the statement follows from the explicit relative Nadel
input (Theorem~\ref{thm:import-sg-ep-char0}; Remark~\ref{rmk:sg-char0}); the case $X=\PP^n$
also has an independent proof
(Theorem~\ref{thm:c0-pn-ep}). Corollary~\ref{cor:sg-smooth-model} transfers the result to the smooth models of rank one of
\S\ref{sec:global} (Theorem~\ref{thm:gl-smooth-ep}), and Remark~\ref{rmk:sg-nef} explains the passage to arbitrary nef classes.

\begin{convention}\label{conv:sg-setting}
\uses{conv:pg-setting, conv:wu-graded, lem:wu-proj-chart, def:pn-gauss, def:pg-subvariety-fs, def:pg-subvariety-psh}
$k$ is algebraically closed with the trivial absolute value. We use Convention~\ref{conv:pg-setting}: $R_M=k[x_0,\dots,x_M]$,
$I\subseteq R_M$ a homogeneous prime not containing all $x_i$, $X=V(I)$, $S=R_M/I$ with the images $\bar x_i$ of the $x_i$,
$\mathfrak m_+=(\bar x_0,\dots,\bar x_M)\subseteq S$, and $\mathcal J=\{i:x_i\notin I\}$. For $i\in\mathcal J$,
$A_i=S_{(\bar x_i)}$ is the chart algebra of Lemma~\ref{lem:wu-proj-chart}. All $A_i$ have the same fraction field
$k(X)$, the field of fractions $\bar f/\bar g$ of forms of equal degree with $g\notin I$, and $n=\operatorname{trdeg}_kk(X)$.
We say that $X$ has \emph{smooth charts} if every $A_i$, $i\in\mathcal J$, is a smooth $k$-algebra.

$T=S[t]$ is bigraded by $\deg\bar x_i=(1,0)$ and $\deg t=(0,1)$. An element of $T$ is \emph{$x$-homogeneous} of
\emph{$x$-degree} $c$ if it lies in $S_c[t]$, and \emph{bihomogeneous} if moreover it is a monomial in $t$; an ideal is
bihomogeneous if it is generated by bihomogeneous elements. For an ideal $\mathfrak b\subseteq T$ its \emph{saturation}
is $\mathfrak b^{\mathrm{sat}}=\{f\in T:\mathfrak m_+^rf\subseteq\mathfrak b\text{ for some }r\}$.

For $v\in X^\an$ the Gauss extension $\sigma(v)$ (Definition~\ref{def:pn-gauss}) vanishes on $I[t]$ and therefore is a
multiplicative seminorm on $T$, with $|\sum_\lambda\bar g_\lambda t^\lambda|_{\sigma(v)}=\max_\lambda|g_\lambda|_ve^{-\lambda}$ and
$|f|_{\sigma(v)}\le1$ for $f\in T$. For a subset $\mathfrak b\subseteq T$ we put
$\phi_{\mathfrak b}(v)=\sup_{f\in\mathfrak b}\log|f|_{\sigma(v)}$. The Fubini--Study functions $\mathcal H(X,a)\supseteq\mathcal H_{\mathrm{fin}}(X,a)$
and the classes $\PSH(X,\theta)$, $\theta\in\R$, are those of Definitions~\ref{def:pg-subvariety-fs}
and~\ref{def:pg-subvariety-psh}.
\end{convention}

\subsection{Envelopes in Fubini--Study systems}\label{subsec:sg-envelopes}

In this subsection $\mathcal F$ is a Fubini--Study system of rank $r$ on a compact Hausdorff space $(K,\star)$
(Definition~\ref{def:pg-fs-system}), with function classes $\mathcal H(a)\supseteq\mathcal H_{\mathrm{fin}}(a)$, $a\in\Q^r_{\ge0}$, and
psh classes $\PSH(\theta)$, $\theta\in\R^r$ (Definition~\ref{def:pg-psh}).

\begin{definition}[Envelopes and the envelope property {\cite[\S2.3]{FHX}, \cite[\S5.2]{BJ22}}]\label{def:sg-ep}
\lean{EnvelopeProperty.FSSystem.envelope, EnvelopeProperty.FSSystem.HasEnvelopeProperty, EnvelopeProperty.FSSystem.IsPL, EnvelopeProperty.ProjSubvariety.envelope, EnvelopeProperty.ProjSubvariety.HasEnvelopeProperty, EnvelopeProperty.ProjSubvariety.IsPL, EnvelopeProperty.ProjSubvariety.isPL_iff}
\leanok
\uses{def:pg-fs-system, def:pg-psh, def:pn-ep}
For $\theta\in\R^r$ and a bounded function $\phi\colon K\to\R$, the \emph{psh envelope} is
$\mathrm P_\theta(\phi)=\sup\{\tau\in\PSH(\theta):\tau\le\phi\}$, a pointwise supremum without regularization (equal to $-\infty$
if the set is empty). The class $\theta$ has the \emph{envelope property} if for every nonempty family $(u_i)$ in
$\PSH(\theta)$ that is bounded above, the upper semicontinuous regularization $(\sup_iu_i)^\star$ lies in $\PSH(\theta)$.
$\operatorname{PL}(K)$ is the set of differences $f-g$ with $f\in\mathcal H_{\mathrm{fin}}(a)$, $g\in\mathcal H_{\mathrm{fin}}(b)$,
$a,b\in\Q^r_{\ge0}$. For $X$ we use the system of Lemma~\ref{lem:pg-subvariety-fs-system} and write $\operatorname{PL}(X)$.
\end{definition}

\begin{lemma}[Envelopes {\cite[\S2.3]{FHX}}]\label{lem:sg-env-basic}
\lean{EnvelopeProperty.FSSystem.envelope_le, EnvelopeProperty.FSSystem.le_envelope, EnvelopeProperty.FSSystem.envelope_mono, EnvelopeProperty.FSSystem.envelope_eq_self, EnvelopeProperty.FSSystem.const_le_envelope, EnvelopeProperty.FSSystem.iInf_le_envelope, EnvelopeProperty.FSSystem.envelope_add_const, EnvelopeProperty.FSSystem.envelope_smul, EnvelopeProperty.FSSystem.HasEnvelopeProperty.isPSH_uscReg_iSup, EnvelopeProperty.FSSystem.HasEnvelopeProperty.isPSH_envelope, EnvelopeProperty.FSSystem.HasEnvelopeProperty.isPSH_envelope_of_continuous, EnvelopeProperty.ProjSubvariety.envelope_add_const, EnvelopeProperty.ProjSubvariety.envelope_mul, EnvelopeProperty.ProjSubvariety.iInf_le_envelope}
\leanok
\uses{def:sg-ep, thm:pg-psh-basic, prop:pg-maxprin}
Let $\theta\in\R^r$, $\theta\ge0$, and $\phi$ bounded. Then $\mathrm P_\theta$ is increasing in $\phi$,
$\inf\phi\le\mathrm P_\theta(\phi)\le\phi$, $\mathrm P_\theta(\phi+c)=\mathrm P_\theta(\phi)+c$ for $c\in\R$, and
$\mathrm P_{t\theta}(t\phi)=t\,\mathrm P_\theta(\phi)$ for real $t>0$.
\end{lemma}

\begin{proof}
Monotonicity and $\mathrm P_\theta(\phi)\le\phi$ are clear. The constant $\inf\phi$ is $\theta$-psh
(Theorem~\ref{thm:pg-psh-basic}(2)), so it is a competitor. By Theorem~\ref{thm:pg-psh-basic}(1),(2), $u\mapsto u+c$ is a
bijection of $\PSH(\theta)$ and $u\mapsto tu$ a bijection $\PSH(\theta)\to\PSH(t\theta)$; both transform the constraint
$u\le\phi$ into $u+c\le\phi+c$, resp.\ $tu\le t\phi$.
\end{proof}

\begin{lemma}[Approximation of classes {\cite[Lemma~5.9]{BJ22}, \cite[Theorem~3.1]{FHX}}]\label{lem:sg-envppty}
\lean{EnvelopeProperty.FSSystem.HasEnvelopeProperty.of_tendsto, EnvelopeProperty.FSSystem.hasEnvelopeProperty_of_forall_rat_gt, EnvelopeProperty.ProjSubvariety.HasEnvelopeProperty.of_tendsto, EnvelopeProperty.ProjSubvariety.hasEnvelopeProperty_of_forall_rat_gt}
\leanok
\uses{def:sg-ep, thm:pg-psh-basic, thm:pg-psh2}
Let $\theta\in\R^r$ be the limit of classes $\theta_j\in\R^r$ with $\theta_j\ge\theta$ coordinatewise, each having the envelope
property. Then $\theta$ has the envelope property.
\end{lemma}

\begin{proof}
Let $(u_i)$ be a nonempty bounded-above family in $\PSH(\theta)$ and $u=(\sup_iu_i)^\star$. By
Theorem~\ref{thm:pg-psh-basic}(3), $u_i\in\PSH(\theta_j)$ for all $j$, so $u\in\PSH(\theta_j)$ for all $j$. The constant
sequence $u$ is decreasing, its classes $\theta_j$ tend to $\theta$, and $u\ge u_{i_0}\not\equiv-\infty$; PSH2
(Theorem~\ref{thm:pg-psh2}(3)) gives $u\in\PSH(\theta)$.
\end{proof}

\begin{lemma}[PL functions {\cite[Lemma~2.3, Theorem~2.2]{BJ22}}]\label{lem:sg-pl}
\lean{EnvelopeProperty.FSSystem.IsFinFS.isPL, EnvelopeProperty.FSSystem.isPL_const, EnvelopeProperty.FSSystem.IsPL.continuous, EnvelopeProperty.FSSystem.IsPL.add, EnvelopeProperty.FSSystem.IsPL.sub, EnvelopeProperty.FSSystem.IsPL.smul, EnvelopeProperty.FSSystem.IsPL.max, EnvelopeProperty.FSSystem.IsPL.min, EnvelopeProperty.FSSystem.plSubmodule, EnvelopeProperty.FSSystem.exists_isPL_abs_sub_lt, EnvelopeProperty.FSSystem.dense_setOf_isPL, EnvelopeProperty.FSSystem.exists_isPL_gt, EnvelopeProperty.FSSystem.directedOn_isPL_gt, EnvelopeProperty.FSSystem.iInf_isPL_gt, EnvelopeProperty.FSSystem.exists_antitone_isPL_net}
\leanok
\uses{def:sg-ep, def:pg-fs-system, lem:bj-pl-lattice, thm:bj-pl-dense, lem:bj-usc-pl}
\begin{enumerate}
\item $\operatorname{PL}(K)$ is a $\Q$-vector space of continuous real functions which contains the rational constants and
is stable under finite maxima and minima.
\item If $\operatorname{PL}(K)$ separates the points of $K$, it is dense in $C^0(K,\R)$ for uniform convergence.
\item In the situation of (2), for every upper semicontinuous $h\colon K\to[-\infty,\infty)$ the set
$\{\psi\in\operatorname{PL}(K):\psi>h\}$ is nonempty, directed downwards, and its pointwise infimum is $h$.
\end{enumerate}
\end{lemma}

\begin{proof}
(1) Functions in $\mathcal H_{\mathrm{fin}}(a)$ are continuous and real valued. Sums of differences are differences of sums,
by (FS4); for $c\in\Q_{>0}$, $c(f-g)=cf-cg$ with $cf\in\mathcal H_{\mathrm{fin}}(ca)$ (FS4), and $-(f-g)=g-f$; a rational constant $c$
lies in $\mathcal H_{\mathrm{fin}}(0)$ (FS3), so $c=c-0$. For maxima,
$\max(f_1-g_1,f_2-g_2)=\max(f_1+g_2,f_2+g_1)-(g_1+g_2)$, where $f_1+g_2\in\mathcal H_{\mathrm{fin}}(a_1+b_2)$ and
$f_2+g_1\in\mathcal H_{\mathrm{fin}}(a_2+b_1)$ both lie in $\mathcal H_{\mathrm{fin}}(c)$ for the coordinatewise maximum $c$ of the two
classes (FS2), hence so does their maximum (FS3). Minima follow by negation.
(2) and (3): the proofs of Theorem~\ref{thm:bj-pl-dense} and Lemma~\ref{lem:bj-usc-pl} use only (1), the separation of
points, the lattice form of the Stone--Weierstrass theorem on a compact Hausdorff space, and Urysohn's lemma; they apply
verbatim to $K$.
\end{proof}

\begin{proposition}[Envelope property via PL obstacles {\cite[Lemma~1.3]{BJadd}, \cite[Proposition~2.5]{FHX}}]\label{prop:sg-envelope-criterion}
\lean{EnvelopeProperty.FSSystem.hasEnvelopeProperty_of_forall_isPL, EnvelopeProperty.FSSystem.hasEnvelopeProperty_iff_forall_isPL, EnvelopeProperty.ProjSubvariety.hasEnvelopeProperty_of_forall_isPL, EnvelopeProperty.ProjSubvariety.hasEnvelopeProperty_iff_forall_isPL}
\leanok
\uses{lem:sg-pl, lem:sg-env-basic, thm:pg-psh2, def:sg-ep}
Assume that $\operatorname{PL}(K)$ separates the points of $K$. Let $\theta\in\R^r$, $\theta\ge0$, and assume that
$\mathrm P_\theta(\psi)\in\PSH(\theta)$ for every $\psi\in\operatorname{PL}(K)$. Then $\theta$ has the envelope property.
\end{proposition}

\begin{proof}
Let $(u_i)$ be a nonempty bounded-above family in $\PSH(\theta)$ and $h=(\sup_iu_i)^\star$. For $\psi\in\operatorname{PL}(K)$ with
$\psi>h$ we have $u_i\le\psi$, hence $u_i\le\mathrm P_\theta(\psi)$ for all $i$; as $\mathrm P_\theta(\psi)$ is psh, hence upper
semicontinuous, $h\le\mathrm P_\theta(\psi)\le\psi$. The set of such $\psi$ is directed downwards and has infimum $h$
(Lemma~\ref{lem:sg-pl}(3)), and $\mathrm P_\theta$ is increasing (Lemma~\ref{lem:sg-env-basic}); so the functions
$\mathrm P_\theta(\psi)$ form a decreasing net of $\theta$-psh functions with limit $h$. Since $h\ge u_{i_0}\not\equiv-\infty$,
PSH2 (Theorem~\ref{thm:pg-psh2}(3)) gives $h\in\PSH(\theta)$.
\end{proof}

\begin{lemma}[Fubini--Study envelopes {\cite[Lemma~1.5]{BJadd}, \cite[Lemma~5.19]{BJ22}}]\label{lem:sg-fs-envelope}
\lean{EnvelopeProperty.FSSystem.envelope_eq_iSup_isFinFS, EnvelopeProperty.FSSystem.envelope_coe_eq_iSup_isFinFS, EnvelopeProperty.ProjSubvariety.envelope_eq_iSup_isFinFS}
\leanok
\uses{thm:pg-regularization, def:sg-ep, def:pg-fs-system}
Let $a\in\Q^r$ have positive coordinates and let $\phi\colon K\to\R$ be bounded and lower semicontinuous. Then
$\mathrm P_a(\phi)=\sup\{\psi\in\mathcal H_{\mathrm{fin}}(a):\psi\le\phi\}$.
\end{lemma}

\begin{proof}
Functions in $\mathcal H_{\mathrm{fin}}(a)$ are $a$-psh (Theorem~\ref{thm:pg-psh2}(1)), which gives $\ge$. Let $\tau\in\PSH(a)$ with
$\tau\le\phi$, put $C=\sup|\phi|+1$ and $\tau'=\max(\tau,-C)$, which is $a$-psh (Theorem~\ref{thm:pg-psh-basic}(2)) and $\le\phi$;
it suffices to show $\tau'\le\mathrm Q$ for the right-hand side $\mathrm Q$. Let $\epsilon\in\Q_{>0}$. By
Theorem~\ref{thm:pg-regularization}(1) there is a decreasing net $\tau_\alpha\in\mathcal H_{\mathrm{fin}}(a_\alpha)$ with limit $\tau'$,
$a_\alpha\ge a$ and $a_\alpha\to a$. The open sets $\{\tau_\alpha<\phi+\epsilon\}$ ($\phi-\tau_\alpha$ is lower semicontinuous) increase
and cover the compact $K$, so $\tau_\alpha<\phi+\epsilon$ for $\alpha\ge\alpha_0$; also $\tau_\alpha\ge\tau'\ge-C$. Choose rationals
$s_\alpha\in(0,1]$ with $s_\alpha a_\alpha\le a$ coordinatewise and $s_\alpha\to1$ (possible as $a$ has positive coordinates). By (FS4),
(FS2) and (FS3), $\psi_\alpha=s_\alpha\tau_\alpha-\epsilon-\delta_\alpha\in\mathcal H_{\mathrm{fin}}(a)$ for rationals $\delta_\alpha\ge(1-s_\alpha)C$, and for
$\alpha\ge\alpha_0$, $s_\alpha\tau_\alpha=\tau_\alpha-(1-s_\alpha)\tau_\alpha\le\tau_\alpha+(1-s_\alpha)C<\phi+\epsilon+\delta_\alpha$, so $\psi_\alpha\le\phi$ and
$\psi_\alpha\le\mathrm Q$. Letting $\alpha\to\infty$ pointwise and choosing $\delta_\alpha\to0$ gives $\tau'-\epsilon\le\mathrm Q$; let
$\epsilon\to0$.
\end{proof}

\begin{theorem}[Reduction to integral classes and nonpositive PL obstacles {\cite[Theorem~3.1]{FHX}, \cite[Lemma~5.9]{BJ22}}]\label{thm:sg-reduction}
\lean{EnvelopeProperty.FSSystem.hasEnvelopeProperty_of_forall_nat, EnvelopeProperty.ProjSubvariety.hasEnvelopeProperty_of_forall_nat, EnvelopeProperty.exists_nat_mul_rat_vec_eq_nat}
\leanok
\uses{lem:sg-envppty, prop:sg-envelope-criterion, lem:sg-env-basic, lem:sg-pl, prop:pg-maxprin, thm:pg-psh-basic}
Assume that $\operatorname{PL}(K)$ separates the points of $K$, and that $\mathrm P_a(\phi)\in\PSH(a)$ for every $a\in\N^r$ with
positive coordinates and every $\phi\in\operatorname{PL}(K)$ with $\phi\le0$. Then every $\theta\in\R^r$ has the envelope property.
\end{theorem}

\begin{proof}
If $\theta\not\ge0$, then $\PSH(\theta)=\emptyset$ by Proposition~\ref{prop:pg-maxprin}, and the property holds vacuously. If
$\theta\ge0$, choose $a_j\in\Q^r$ with positive coordinates, $a_j\ge\theta$ and $a_j\to\theta$; by Lemma~\ref{lem:sg-envppty} it
suffices to treat $\theta=a\in\Q^r$ with positive coordinates, and by Proposition~\ref{prop:sg-envelope-criterion} to show
$\mathrm P_a(\psi)\in\PSH(a)$ for $\psi\in\operatorname{PL}(K)$. Choose $m\in\mathbf{Z}_{>0}$ with $ma\in\N^r$ and a rational $c\ge\max m\psi$. Then
$m\psi-c\in\operatorname{PL}(K)$ is $\le0$ (Lemma~\ref{lem:sg-pl}(1)), so $\mathrm P_{ma}(m\psi-c)\in\PSH(ma)$ by hypothesis, and
$\mathrm P_a(\psi)=m^{-1}(\mathrm P_{ma}(m\psi-c)+c)$ (Lemma~\ref{lem:sg-env-basic}) lies in $\PSH(a)$ by
Theorem~\ref{thm:pg-psh-basic}(1),(2).
\end{proof}

\subsection{Divisorial points and divisorial comparison on \texorpdfstring{$X$}{X}}\label{subsec:sg-comparison}

\begin{lemma}[The punctured cone]\label{lem:sg-cone-charts}
\lean{EnvelopeProperty.ProjSubvariety.CoordRing, EnvelopeProperty.ProjSubvariety.ConeRing, EnvelopeProperty.ProjSubvariety.coneChartEquiv, EnvelopeProperty.ProjSubvariety.chartPolyToFrac_injective, EnvelopeProperty.ProjSubvariety.gaussToFrac_injective, EnvelopeProperty.ProjSubvariety.trdeg_fractionRing_coneRing, EnvelopeProperty.ProjSubvariety.HasSmoothCharts, EnvelopeProperty.ProjSubvariety.smooth_coneAway, EnvelopeProperty.ProjSubvariety.isBihom_iff_chartModel, EnvelopeProperty.ProjSubvariety.modelToFrac_injective, EnvelopeProperty.ProjSubvariety.isFractionRing_chartModel, EnvelopeProperty.ProjSubvariety.dim, EnvelopeProperty.ProjSubvariety.trdeg_chartField_eq_dim}
\leanok
\uses{conv:sg-setting, lem:wu-proj-chart}
Let $i\in\mathcal J$ and let $Y$, $t$ be indeterminates.
\begin{enumerate}
\item The $k$-algebra homomorphism $A_i[Y,Y^{-1}][t]\to T_{\bar x_i}$ which is the inclusion on $A_i$ and maps $Y\mapsto\bar x_i$,
$t\mapsto t$ is an isomorphism; it matches the bigrading of $T_{\bar x_i}$ with $\deg Y=(1,0)$, $\deg t=(0,1)$ and $A_i$ in
bidegree $(0,0)$.
\item $T$ is a domain, $\operatorname{Frac}T=k(X)(\bar x_i,t)$ with $\bar x_i,t$ algebraically independent over $k(X)$, and
$\operatorname{trdeg}_k\operatorname{Frac}T=n+2$.
\item If $A_i$ is smooth over $k$, so is $T_{\bar x_i}$.
\end{enumerate}
\end{lemma}

\begin{proof}
(1) An element of $S_{\bar x_i}$ is $\bar f/\bar x_i^N$ with $f=\sum_cf_c$, $f_c$ homogeneous of degree $c$, and
$\bar f/\bar x_i^N=\sum_c(\bar f_c/\bar x_i^c)\,\bar x_i^{c-N}$ with $\bar f_c/\bar x_i^c\in A_i$; so $A_i[Y^{\pm1}]\to S_{\bar x_i}$ is surjective.
If $\sum_ca_c\bar x_i^c=0$ with $a_c\in A_i$, each summand is the component of degree $c$ of the $\mathbf{Z}$-graded ring
$S_{\bar x_i}$, hence vanishes, and $a_c=0$ because $\bar x_i$ is a unit. Adjoining $t$ gives the claim.
(2) $S$ is a domain since $I$ is prime, and $T\subseteq T_{\bar x_i}\cong A_i[Y^{\pm1}][t]$, whose fraction field is
$k(X)(Y,t)$ with $Y,t$ algebraically independent over $\operatorname{Frac}A_i=k(X)$. (3) Polynomial algebras, localizations and
composites of smooth algebras are smooth.
\end{proof}

\begin{lemma}[Functions attached to ideals of $T$ {\cite[\S2.1]{FHX}}]\label{lem:sg-phi-basic}
\lean{EnvelopeProperty.ProjSubvariety.phiIdeal_sup_ker_coneMk, EnvelopeProperty.ProjSubvariety.phiIdeal_comap_coneIrrelevant, EnvelopeProperty.ProjSubvariety.phiIdeal_comap_saturation, EnvelopeProperty.ProjSubvariety.isFS_inv_mul_phiIdeal, EnvelopeProperty.ProjSubvariety.gaussExt_val_eq_zero_of_coneMk_eq_zero, EnvelopeProperty.ProjSubvariety.logAbs_gaussExt_eq_of_coneMk_eq}
\leanok
\uses{conv:sg-setting, def:pg-subvariety-fs, lem:pg-subvariety-fs-system}
Let $\mathfrak b,\mathfrak b'\subseteq T$ be ideals.
\begin{enumerate}
\item $\phi_{\mathfrak b}\le0$; $\phi_{\mathfrak b}\le\phi_{\mathfrak b'}$ if $\mathfrak b\subseteq\mathfrak b'$; $\phi_{\mathfrak b}=\max_k\log|g_k|_{\sigma(\cdot)}$ for
finitely many generators $g_k$ of $\mathfrak b$; $\phi_{\mathfrak b\mathfrak b'}=\phi_{\mathfrak b}+\phi_{\mathfrak b'}$.
\item $\phi_{\mathfrak b^{\mathrm{sat}}}=\phi_{\mathfrak b}$.
\item If $\mathfrak b\ne0$ is generated by finitely many bihomogeneous elements of $x$-degree $\le D$, then
$\phi_{\mathfrak b}\in\mathcal H(X,D)$, with integer weights.
\end{enumerate}
\end{lemma}

\begin{proof}
(1) Elements of $T$ have $|\cdot|_{\sigma(v)}\le1$; for $f=\sum_ka_kg_k$, $|f|\le\max_k|a_k||g_k|\le\max_k|g_k|$ by the ultrametric
inequality, and the $g_k$ belong to $\mathfrak b$. The products $g_kg'_l$ generate $\mathfrak b\mathfrak b'$, and $\sigma(v)$ is multiplicative.
(2) $\mathfrak b^{\mathrm{sat}}$ is finitely generated, so $\mathfrak m_+^r\mathfrak b^{\mathrm{sat}}\subseteq\mathfrak b\subseteq\mathfrak b^{\mathrm{sat}}$ for some $r$,
and $\phi_{\mathfrak m_+^r}(v)=r\max_i\log|x_i|_v=0$ since $v$ is normalized; apply (1).
(3) Drop the zero generators and write the others as $\bar s_kt^{\lambda_k}$ with forms $s_k\notin I$ of degree $D_k\le D$. As
$\max_{|\alpha|=D-D_k}\log|x^\alpha|_v=0$,
$\phi_{\mathfrak b}(v)=\max_k(\log|s_k|_v-\lambda_k)=\max_{k,\,|\alpha|=D-D_k}(\log|s_kx^\alpha|_v-\lambda_k)$, a Fubini--Study function of class
$D$ with forms of degree $D$, among which $s_kx_j^{D-D_k}\notin I$ for $j\in\mathcal J$ ($I$ prime).
\end{proof}

\begin{lemma}[PL functions separate points of $X^\an$ {\cite[Lemma~2.3]{BJ22}}]\label{lem:sg-pl-separate}
\lean{EnvelopeProperty.ProjSubvariety.exists_isFiniteFS_ne, EnvelopeProperty.ProjSubvariety.separatesPoints_isPL}
\leanok
\uses{lem:pg-subvariety-fs-system, def:sg-ep, def:pn-an}
$\operatorname{PL}(X)$ separates the points of $X^\an$.
\end{lemma}

\begin{proof}
Let $v\ne w$ in $X^\an$. Both are graded seminorms on $R_M$, so they differ on some form $s$, of degree $c\ge1$ (they agree
on constants), and $s\notin I$ (both vanish on $I$). For $N\in\N$ the function
$f_N=\max\{\log|s|,\max_i\log|x_i^c|-N\}$ lies in $\mathcal H_{\mathrm{fin}}(X,c)$: the weights are rational, some $x_i^c\notin I$, and
the forms have no common zero on $X^\an$ because $\max_i|x_i|=1$ there. On $X^\an$, $f_N=\max\{\log|s|,-N\}$. If
$|s|_v,|s|_w>0$, then $f_N(v)=\log|s|_v\ne\log|s|_w=f_N(w)$ for large $N$; if, say, $|s|_v=0<|s|_w$, then $f_N(v)=-N\ne f_N(w)$
for large $N$. So $f_N=f_N-0\in\operatorname{PL}(X)$ separates $v$ and $w$.
\end{proof}

\begin{definition}[Divisorial points of $X$ {\cite[\S2.1]{FHX}, \cite[\S1.3]{BJ22}}]\label{def:sg-div}
\lean{EnvelopeProperty.ProjSubvariety.IsDivisorial, EnvelopeProperty.ProjSubvariety.isDivisorial_of_val_eq}
\leanok
\uses{conv:sg-setting, def:pn-div}
A point $v\in X^\an$ is \emph{divisorial} if there are $c\in\Q_{>0}$ and a discrete valuation ring $R'$ with
$T\subseteq R'\subseteq\operatorname{Frac}T$, essentially of finite type over $k$, such that $\log|f|_{\sigma(v)}=-c\operatorname{ord}_{R'}(f)$ for all
$f\in T$. We write $X^{\divv}$ for the set of divisorial points. For such $v$ and an ideal $\mathfrak b\subseteq T$,
$\phi_{\mathfrak b}(v)=-c\operatorname{ord}_{R'}(\mathfrak b)$, and $\bar x_j$ is a unit of $R'$ for every $j$ with $|x_j|_v=1$.
\end{definition}

\begin{lemma}[Examples of divisorial points]\label{lem:sg-div-points}
\lean{EnvelopeProperty.ProjSubvariety.isDivisorial_triv, EnvelopeProperty.ProjSubvariety.chartDivPoint, EnvelopeProperty.ProjSubvariety.isDivisorial_chartDivPoint, EnvelopeProperty.ProjSubvariety.chartDVR, EnvelopeProperty.ProjSubvariety.chartDVR_ord, EnvelopeProperty.ProjSubvariety.tAdicDVR, EnvelopeProperty.ProjSubvariety.tAdicDVR_ord}
\leanok
\uses{def:sg-div, lem:sg-cone-charts, lem:bj-graded-extension, lem:bj-gauss-dvr}
\begin{enumerate}
\item The trivial point $\star$ is divisorial.
\item Let $j\in\mathcal J$, $B_j=A_j[t]$, and let $V$ be a discrete valuation ring with $B_j\subseteq V\subseteq\operatorname{Frac}B_j$, essentially
of finite type over $k$, with $b\coloneqq\operatorname{ord}_V(t)>0$ and $\operatorname{ord}_V(\sum_\lambda f_\lambda t^\lambda)=\min_\lambda(\operatorname{ord}_V(f_\lambda)+\lambda b)$ for
$f_\lambda\in A_j$. There is a divisorial point $v_V\in X^\an$ with $|x_j|_{v_V}=1$ and
$\log|f|_{\sigma(v_V)}=-\operatorname{ord}_V(f/\bar x_j^c)/b$ for every $f\in T$ which is $x$-homogeneous of degree $c$.
\end{enumerate}
\end{lemma}

\begin{proof}
(1) $|\sum_\lambda\bar g_\lambda t^\lambda|_{\sigma(\star)}=\max\{e^{-\lambda}:\bar g_\lambda\ne0\}=e^{-\operatorname{ord}_t}$. The ideal $tT$ is prime
($T/tT=S$), and $R'=T_{(t)}$ is a Noetherian local domain with nonzero principal maximal ideal, hence a discrete valuation
ring, essentially of finite type, with $\operatorname{ord}_{R'}=\operatorname{ord}_t$; take $c=1$.
(2) For a form $s$ of degree $c$ put $N(s)=\exp(-\operatorname{ord}_V(\bar s/\bar x_j^c)/b)$ (so $N(s)=0$ for $s\in I$). $N$ is multiplicative,
ultrametric on forms of equal degree, equal to $1$ on $k^\times$, with $N(x_j)=1$ and $N(x_i)\le1$ because
$\bar x_i/\bar x_j\in A_j\subseteq V$. Lemma~\ref{lem:bj-graded-extension} gives a point $v_V$ of $(\PP^M)^\an$ with $|s|=N(s)$ on forms; it
vanishes on the homogeneous elements of $I$, hence on $I$, so $v_V\in X^\an$. For $f=\sum_\lambda\bar s_\lambda t^\lambda$ $x$-homogeneous of degree
$c$, $\log|f|_{\sigma(v_V)}=\max_\lambda(\log N(s_\lambda)-\lambda)=-b^{-1}\min_\lambda(\operatorname{ord}_V(\bar s_\lambda/\bar x_j^c)+\lambda b)=-b^{-1}\operatorname{ord}_V(f/\bar x_j^c)$.
By Lemma~\ref{lem:sg-cone-charts}(2), $\operatorname{Frac}T=(\operatorname{Frac}B_j)(\bar x_j)$ with $\bar x_j$ transcendental over $\operatorname{Frac}B_j$, so by
Lemma~\ref{lem:bj-gauss-dvr} $R'=V[\bar x_j]_{\mathfrak m_VV[\bar x_j]}\subseteq\operatorname{Frac}T$ is a discrete valuation ring, essentially of finite type,
with $\operatorname{ord}_{R'}(\sum_ia_i\bar x_j^i)=\min_i\operatorname{ord}_V(a_i)$. It contains $\bar x_i=(\bar x_i/\bar x_j)\bar x_j$ and $t$, hence $T$. For $f\in T$ with
$x$-homogeneous components $f_c$, $f=\sum_c(f_c/\bar x_j^c)\bar x_j^c$ with $f_c/\bar x_j^c\in B_j$, so
$\operatorname{ord}_{R'}(f)=\min_c\operatorname{ord}_V(f_c/\bar x_j^c)=-b\max_c\log|f_c|_{\sigma(v_V)}=-b\log|f|_{\sigma(v_V)}$, the last equality because $v_V$ is graded.
Take $c=1/b$.
\end{proof}

\begin{lemma}[Invariance under the torus {\cite[Lemma~2.13]{BJ22}}]\label{lem:sg-rees-invariant}
\lean{EnvelopeProperty.ProjSubvariety.reesDVR, EnvelopeProperty.ProjSubvariety.reesDVR_torus, EnvelopeProperty.ProjSubvariety.reesOrdX_pos, EnvelopeProperty.ProjSubvariety.isDivisorial_reesPoint}
\leanok
\uses{def:bj-rees, lem:bj-rees-invariant}
Let $A$ be a finitely generated $k$-algebra which is a domain, $B=A[t]$, and $\mathfrak i=(g_1,\dots,g_s)\subseteq B$ with $g_\ell=h_\ell t^{e_\ell}$,
$0\ne h_\ell\in A$, and $t^N\in\mathfrak i$ for some $N$. For every $V\in\mathrm{RV}(\mathfrak i;g)$ (Definition~\ref{def:bj-rees}),
$\operatorname{ord}_V(t)>0$ and $\operatorname{ord}_V(\sum_\lambda f_\lambda t^\lambda)=\min_\lambda(\operatorname{ord}_V(f_\lambda)+\lambda\operatorname{ord}_V(t))$ for $f_\lambda\in A$.
\end{lemma}

\begin{proof}
Let $V=(E_\ell)_P$ with $D_\ell=B[g/g_\ell]$, $E_\ell$ its integral closure and $P$ a minimal prime of $g_\ell E_\ell$. Since $t^N\in\mathfrak i$,
$t^N\in g_\ell D_\ell\subseteq P$, so $t\in P$ and $\operatorname{ord}_V(t)>0$. For $\lambda\in k^\times$ the $A$-automorphism $\alpha_\lambda\colon t\mapsto\lambda t$ of
$B$ extends to $\operatorname{Frac}B$, maps $g_i/g_\ell$ to $\lambda^{e_i-e_\ell}g_i/g_\ell$, hence preserves $D_\ell$, $E_\ell$ and $g_\ell E_\ell$, and permutes the
finite set of minimal primes of $g_\ell E_\ell$. The stabilizer of $P$ is a subgroup of finite index of the divisible group
$k^\times$, hence equal to $k^\times$, so $\operatorname{ord}_V\circ\alpha_\lambda=\operatorname{ord}_V$. For $f=\sum_{\lambda\le L}f_\lambda t^\lambda$ choose distinct
$\mu_0,\dots,\mu_L\in k^\times$; inverting the Vandermonde matrix writes each $f_\lambda t^\lambda$ as a $k$-linear combination of the
$\alpha_{\mu_i}(f)$, so $\operatorname{ord}_V(f_\lambda t^\lambda)\ge\operatorname{ord}_V(f)$; the reverse inequality is the ultrametric inequality.
\end{proof}

\begin{lemma}[Finite Fubini--Study functions as ideals {\cite[Proposition~2.25]{BJ22}}]\label{lem:sg-finite-fs}
\lean{EnvelopeProperty.ProjSubvariety.IsFinFS.exists_coneSup}
\leanok
\uses{lem:bj-graded-extension, conv:sg-setting, lem:sg-phi-basic}
Let $\psi=m^{-1}\max_i(\log|s_i|+\lambda_i)\in\mathcal H_{\mathrm{fin}}(X,b)$ with forms $s_i\notin I$ of degree $mb$, $\lambda_i\in\mathbf{Z}_{\le0}$ and
$\max_i\lambda_i=0$. Then $\mathfrak m_+^e\subseteq(\bar s_i)_iS$ for some $e\ge mb$; the ideal $\mathfrak i\subseteq T$ generated by the
$g_i=\bar s_it^{-\lambda_i}$ satisfies $\phi_{\mathfrak i}=m\psi$ on $X^\an$; and for $N=\max_i(-\lambda_i)$ and every $j\in\mathcal J$,
$t^N$ lies in the ideal $\mathfrak i^{(j)}$ of $B_j=A_j[t]$ generated by the $g_i^{(j)}=g_i/\bar x_j^{mb}$.
\end{lemma}

\begin{proof}
Let $\mathfrak p\subseteq S$ be a homogeneous prime containing all $\bar s_i$ but not $\mathfrak m_+$. The function $N(f)=0$ for $\bar f\in\mathfrak p$ and
$N(f)=1$ otherwise, on forms $f$, satisfies the hypotheses of Lemma~\ref{lem:bj-graded-extension} (some $\bar x_j\notin\mathfrak p$), and the
resulting point lies in $X^\an$ and has $|s_i|=0$ for all $i$, contradicting $\psi\in\mathcal H_{\mathrm{fin}}(X,b)$. Hence every homogeneous
prime containing $(\bar s_i)$ contains $\mathfrak m_+$; the radical of a homogeneous ideal is the intersection of the homogeneous primes
containing it, so $\mathfrak m_+^e\subseteq(\bar s_i)$ for some $e$, which we may take $\ge mb$. Then
$\phi_{\mathfrak i}=\max_i\log|g_i|_{\sigma(\cdot)}=\max_i(\log|s_i|+\lambda_i)=m\psi$ (Lemma~\ref{lem:sg-phi-basic}(1)). Write
$\bar x_j^e=\sum_ic_i\bar s_i$ with $c_i\in S_{e-mb}$; then $\bar x_j^et^N=\sum_ic_it^{N+\lambda_i}g_i$, and dividing by $\bar x_j^e$ gives
$t^N=\sum_i(c_i/\bar x_j^{e-mb})t^{N+\lambda_i}g_i^{(j)}$ in $B_j$.
\end{proof}

\begin{lemma}[Shilov boundary on $X$ {\cite[Lemma~2.13]{BJ22}}]\label{lem:sg-shilov}
\lean{EnvelopeProperty.ProjSubvariety.IsFinFS.exists_finset_iSup_sub_eq, EnvelopeProperty.ProjSubvariety.IsFinFS.exists_finset_le_of_le}
\leanok
\uses{lem:sg-finite-fs, thm:bj-rees, def:bj-rees, thm:bj-normalization-finite, lem:sg-rees-invariant, lem:sg-div-points, lem:sg-phi-basic}
For every $\psi\in\mathcal H_{\mathrm{fin}}(X,b)$ there is a finite set $\Sigma\subseteq X^{\divv}$ containing $\star$ such that for every $a\in\Q_{\ge0}$ and
every $\varphi\in\mathcal H(X,a)$ with rational weights, $\sup_{X^\an}(\varphi-\psi)=\max_\Sigma(\varphi-\psi)$.
\end{lemma}

\begin{proof}
Subtracting a rational constant from $\psi$ and replacing $m$ by a multiple and the $s_i$ by powers, we may assume $\psi$ as in
Lemma~\ref{lem:sg-finite-fs}; let $g_i$, $\mathfrak i$, $N$, $B_j$, $\mathfrak i^{(j)}$ be as there. For $j\in\mathcal J$, $B_j$ is a finitely generated domain,
$g_i^{(j)}=(\bar s_i/\bar x_j^{mb})t^{-\lambda_i}$ with $\bar s_i/\bar x_j^{mb}\ne0$, and $t^N\in\mathfrak i^{(j)}$; by Definition~\ref{def:bj-rees},
Lemma~\ref{lem:sg-rees-invariant} and Lemma~\ref{lem:sg-div-points}(2), every $V\in\mathrm{RV}(\mathfrak i^{(j)};g^{(j)})$ yields a divisorial point
$v_V$. Let $\Sigma$ consist of $\star$ and these finitely many points. Write $\varphi=M^{-1}\max_\ell(\log|\sigma_\ell|+\mu_\ell)$ with forms
$\sigma_\ell\notin I$ of degree $Ma$ (terms with $\sigma_\ell\in I$ are identically $-\infty$ on $X^\an$ and are dropped) and $\mu_\ell\in\mathbf{Z}$. Put
$c=\max_\Sigma(\varphi-\psi)$; since $\psi(\star)=0$ and $\varphi(\star)=M^{-1}\max_\ell\mu_\ell$ are finite and the values on $\Sigma$ are rational or
$-\infty$, $c\in\Q$ and $c\ge\mu_\ell/M$ for all $\ell$. Choose $r\ge1$ with $rMc\in\mathbf{Z}$ and put $a_\ell=m(rMc-r\mu_\ell)\in\N$ and
$F_\ell=\bar\sigma_\ell^{rm}t^{a_\ell}\in T$. Multiplying by $rMm$ shows that at any $v$, $(\varphi-\psi)(v)\le c$ if and only if
$\log|F_\ell|_{\sigma(v)}\le rM\phi_{\mathfrak i}(v)=\phi_{\mathfrak i^{rM}}(v)$ for all $\ell$. At $v=v_V\in\Sigma$ (chart $j$),
$\log|F_\ell|_{\sigma(v_V)}=-\operatorname{ord}_V(F_\ell^{(j)})/\operatorname{ord}_V(t)$ and $\phi_{\mathfrak i^{rM}}(v_V)=-rM\operatorname{ord}_V(\mathfrak i^{(j)})/\operatorname{ord}_V(t)$
(Lemma~\ref{lem:sg-div-points}(2), with $F^{(j)}=F/\bar x_j^{\deg_xF}$), so $\operatorname{ord}_V(F_\ell^{(j)})\ge rM\operatorname{ord}_V(\mathfrak i^{(j)})$ for every
$V\in\mathrm{RV}(\mathfrak i^{(j)};g^{(j)})$. Now let $v\in X^\an$ and choose $j$ with $|x_j|_v=1$. For $f\in B_j$ choose $F\in T$ $x$-homogeneous of some degree
$c'$ with $F/\bar x_j^{c'}=f$ and put $|f|_w=|F|_{\sigma(v)}$; this does not depend on $c'$ because $|x_j|_v=1$, and $w$ is a multiplicative
seminorm on $B_j$ with $|\cdot|_w\le1$ on $B_j$ and $|t|_w=e^{-1}$, so $\max_i|g_i^{(j)}|_w\ge|t^N|_w>0$. Theorem~\ref{thm:bj-rees} gives
$|F_\ell^{(j)}|_w\le\max_i|g_i^{(j)}|_w^{rM}$, that is, $\log|F_\ell|_{\sigma(v)}\le\phi_{\mathfrak i^{rM}}(v)$.
\end{proof}

\begin{lemma}[Suprema against PL functions {\cite[Lemma~4.26]{BJ22}}]\label{lem:sg-sup-pl}
\lean{EnvelopeProperty.ProjSubvariety.IsPL.exists_finset_iSup_sub_eq, EnvelopeProperty.ProjSubvariety.IsPL.exists_finset_le_of_le}
\leanok
\uses{lem:sg-shilov, thm:pg-regularization, lem:pg-subvariety-fs-system}
For every $\psi\in\operatorname{PL}(X)$ there is a finite set $\Sigma\subseteq X^{\divv}$ such that $\sup_{X^\an}(u-\psi)=\max_\Sigma(u-\psi)$ for every
$\theta\in\R$ and every $u\in\PSH(X,\theta)$.
\end{lemma}

\begin{proof}
Write $\psi=\psi_1-\psi_2$ with $\psi_\nu\in\mathcal H_{\mathrm{fin}}(X,b_\nu)$ and take $\Sigma$ for $\psi_1$ from Lemma~\ref{lem:sg-shilov}. Let
$u_\alpha\in\mathcal H_{\mathrm{fin}}(X,a_\alpha)$ decrease to $u$ (Theorem~\ref{thm:pg-regularization}(1)). Then
$u_\alpha+\psi_2\in\mathcal H(X,a_\alpha+b_2)$ has rational weights, so for every $v$,
$(u-\psi)(v)\le(u_\alpha+\psi_2-\psi_1)(v)\le\max_\Sigma(u_\alpha-\psi)$. As $\Sigma$ is finite and $u_\alpha$ decreases to $u$ pointwise, the right
side tends to $\max_\Sigma(u-\psi)$.
\end{proof}

\begin{theorem}[Divisorial comparison on $X$ {\cite[Theorem~4.22]{BJ22}, \cite[Proposition~2.1]{FHX}}]\label{thm:sg-comparison}
\lean{EnvelopeProperty.ProjSubvariety.IsPSH.le_of_forall_isDivisorial_le, EnvelopeProperty.ProjSubvariety.IsPSH.le_of_forall_isDivisorial_le_isPL, EnvelopeProperty.ProjSubvariety.IsPSH.eq_of_forall_isDivisorial_eq}
\leanok
\uses{lem:sg-sup-pl, lem:sg-pl, lem:sg-pl-separate}
Let $u\in\PSH(X,\theta)$ and $h\colon X^\an\to[-\infty,\infty)$ upper semicontinuous. If $u\le h$ on $X^{\divv}$, then $u\le h$ on $X^\an$.
\end{theorem}

\begin{proof}
Let $\psi\in\operatorname{PL}(X)$ with $\psi>h$. By Lemma~\ref{lem:sg-sup-pl}, $\sup(u-\psi)=\max_\Sigma(u-\psi)\le0$ since $\Sigma\subseteq X^{\divv}$. So
$u\le\psi$ for all such $\psi$, and $u\le h$ by Lemma~\ref{lem:sg-pl}(3), which applies by Lemma~\ref{lem:sg-pl-separate}.
\end{proof}

\begin{remark}[Relation to the general models]\label{rmk:sg-general-models}
\uses{def:sg-ep, lem:sg-envppty, lem:sg-pl, prop:sg-envelope-criterion, lem:sg-fs-envelope, def:sg-div, lem:sg-div-points, lem:sg-rees-invariant, lem:sg-shilov, thm:sg-comparison, rmk:gl-bridges, def:gl-div, lem:gl-rees-points, lem:gl-shilov, lem:gl-pl, thm:gl-comparison, prop:gl-pl-envelope, lem:gl-fs-envelope, lem:gl-ep-basic}
For the quotient model of $V(I)$ (Remark~\ref{rmk:gl-bridges}), Definition~\ref{def:sg-div}, Lemmas~\ref{lem:sg-div-points}(2),
\ref{lem:sg-rees-invariant}, \ref{lem:sg-shilov}, \ref{lem:sg-pl}, Theorem~\ref{thm:sg-comparison}, Proposition~\ref{prop:sg-envelope-criterion},
Lemmas~\ref{lem:sg-fs-envelope} and~\ref{lem:sg-envppty} are the cases $r=1$ of Definition~\ref{def:gl-div}, Lemmas~\ref{lem:gl-rees-points},
\ref{lem:gl-shilov}, \ref{lem:gl-pl}, Theorem~\ref{thm:gl-comparison}, Proposition~\ref{prop:gl-pl-envelope} and Lemmas~\ref{lem:gl-fs-envelope},
\ref{lem:gl-ep-basic}(1) for the models of \S\ref{sec:global}, with the same proofs. Only one of the two families needs to be formalized; the
other then follows through Remark~\ref{rmk:gl-bridges}.
\end{remark}

\subsection{The ideals \texorpdfstring{$J_m$}{Jm} on \texorpdfstring{$X$}{X}}\label{subsec:sg-Jm}

\begin{definition}[The ideals $J_m$ {\cite[\S3]{FHX}}]\label{def:sg-Jm}
\lean{EnvelopeProperty.ProjSubvariety.obstacleSpace, EnvelopeProperty.ProjSubvariety.obstacleIdeal, EnvelopeProperty.ProjSubvariety.coneObstacleIdeal, EnvelopeProperty.ProjSubvariety.mem_obstacleSpace}
\leanok
\uses{conv:sg-setting, def:bj-Jm}
Let $d\in\mathbf{Z}_{>0}$ and $\phi\colon X^\an\to\R$. For $m\ge1$ let $V_m(\phi)\subseteq S_{md}[t]$ be the set of $x$-homogeneous $g$ of degree
$md$ with $\log|g|_{\sigma(v)}\le m\phi(v)$ for all $v\in X^\an$, and $J_m(\phi)\subseteq T$ the ideal generated by $V_m(\phi)$.
\end{definition}

\begin{lemma}[Properties of $J_m$]\label{lem:sg-Jm-basic}
\lean{EnvelopeProperty.ProjSubvariety.map_mul_mem_obstacleSpace, EnvelopeProperty.ProjSubvariety.monomial_coeff_mem_obstacleSpace, EnvelopeProperty.ProjSubvariety.obstacleIdeal_mul_le, EnvelopeProperty.ProjSubvariety.coneObstacleIdeal_mul_le, EnvelopeProperty.ProjSubvariety.exists_coneObstacleIdeal_eq_span, EnvelopeProperty.ProjSubvariety.coneObstacleIdeal_ne_bot, EnvelopeProperty.ProjSubvariety.isFinFS_obstacleIdeal, EnvelopeProperty.ProjSubvariety.inv_mul_phiIdeal_obstacleIdeal_le, EnvelopeProperty.ProjSubvariety.phiIdeal_comap_coneObstacleIdeal}
\leanok
\uses{def:sg-Jm, lem:sg-phi-basic, lem:pg-subvariety-fs-system}
Let $\phi$ be bounded. Then $V_m(\phi)$ is a $k[t]$-submodule containing the $t$-homogeneous components of its elements;
$J_mJ_r\subseteq J_{m+r}$; $J_m$ is generated by bihomogeneous elements of $x$-degree $md$; $J_m\ne0$; and $m^{-1}\phi_{J_m}\in\mathcal H_{\mathrm{fin}}(X,d)$
with $m^{-1}\phi_{J_m}\le\phi$.
\end{lemma}

\begin{proof}
The ultrametric inequality, $|t|_{\sigma(v)}=e^{-1}$ and the Gauss formula give the first claim, and $V_mV_r\subseteq V_{m+r}$ by
multiplicativity, hence $J_mJ_r\subseteq J_{m+r}$; the $t$-homogeneous components of elements of $V_m$ are bihomogeneous elements of
$V_m$ which generate $J_m$. If $\phi\ge-C$ then $\bar x_j^{md}t^{\lceil mC\rceil}\in V_m$ for $j\in\mathcal J$, so $J_m\ne0$ and
$\phi_{J_m}\ge\max_j\log|x_j^{md}|-\lceil mC\rceil=-\lceil mC\rceil$ on $X^\an$. Choose finitely many bihomogeneous generators of $J_m$ in $V_m$,
including the $\bar x_j^{md}t^{\lceil mC\rceil}$; by Lemma~\ref{lem:sg-phi-basic}(1),(3), $\phi_{J_m}$ is the maximum of their $\log|\cdot|_{\sigma(\cdot)}$, a
Fubini--Study function of class $md$ with integer weights and no common zero on $X^\an$ (it is $\ge-\lceil mC\rceil$), hence in
$\mathcal H_{\mathrm{fin}}(X,md)$, and $\phi_{J_m}\le m\phi$ since each generator lies in $V_m$. Scale by $m^{-1}$ (FS4).
\end{proof}

\begin{proposition}[Envelopes of PL obstacles {\cite[Proposition~2.5]{FHX}, \cite[Lemma~5.22]{BJ22}}]\label{prop:sg-envelope-Jm}
\lean{EnvelopeProperty.ProjSubvariety.envelope_eq_iSup_obstacleIdeal, EnvelopeProperty.ProjSubvariety.envelope_eq_iSup_obstacleIdeal_of_nonpos, EnvelopeProperty.ProjSubvariety.envelope_eq_iSup_coneObstacleIdeal, EnvelopeProperty.ProjSubvariety.inv_mul_phiIdeal_obstacleIdeal_le_mul}
\leanok
\uses{lem:sg-fs-envelope, lem:sg-Jm-basic, def:sg-Jm}
Let $d\in\mathbf{Z}_{>0}$ and $\phi\colon X^\an\to\R$ bounded, lower semicontinuous, with $\phi\le0$. Then
$\mathrm P_d(\phi)=\sup_{m\ge1}m^{-1}\phi_{J_m(\phi)}$, and $m^{-1}\phi_{J_m}\le(mr)^{-1}\phi_{J_{mr}}$ for $r\ge1$.
\end{proposition}

\begin{proof}
$J_m^r\subseteq J_{mr}$ and $\phi_{\mathfrak b^r}=r\phi_{\mathfrak b}$ give the monotonicity. The functions $m^{-1}\phi_{J_m}$ are Fubini--Study functions of
class $d$ with $m^{-1}\phi_{J_m}\le\phi$, hence $\le\mathrm P_d(\phi)$. Conversely, by Lemma~\ref{lem:sg-fs-envelope} it suffices to show
$\psi\le\sup_mm^{-1}\phi_{J_m}$ for $\psi\in\mathcal H_{\mathrm{fin}}(X,d)$ with $\psi\le\phi$. Write $\psi=M^{-1}\max_\ell(\log|s_\ell|+\lambda_\ell)$ with forms
$s_\ell$ of degree $Md$ and, replacing $M$ by a multiple and the $s_\ell$ by powers, $\lambda_\ell\in\mathbf{Z}$; drop the $s_\ell\in I$. Since
$|s_\ell|_\star=1$, evaluation at $\star$ gives $\lambda_\ell\le M\psi(\star)\le M\phi(\star)\le0$. So $g_\ell=\bar s_\ell t^{-\lambda_\ell}\in S_{Md}[t]$ and
$\log|g_\ell|_{\sigma(v)}=\log|s_\ell|_v+\lambda_\ell\le M\psi(v)\le M\phi(v)$; thus $g_\ell\in V_M(\phi)$ and $\psi\le M^{-1}\phi_{J_M}$.
\end{proof}

\subsection{Frobenius roots and test ideals in rings with a \texorpdfstring{$p$}{p}-basis}\label{subsec:sg-pbasis}

From now on $\operatorname{char}k=p>0$; $q$ and $Q$ denote powers of $p$.

\begin{convention}\label{conv:sg-pbasis}
\lean{EnvelopeProperty.Cartier.IsPBasis, EnvelopeProperty.Cartier.IsPBasis.traceIter, EnvelopeProperty.Cartier.IsPBasis.jac}
\leanok
\uses{lem:frobenius-roots, lem:pn-divisorial-trace}
Let $A$ be a commutative ring of characteristic $p$ and $u=(u_1,\dots,u_N)$ a finite family in $A$. Then $u$ is a \emph{$p$-basis}
of $A$ if every $f\in A$ can be written uniquely as $f=\sum_{a\in[0,p)^N}f_a^pu^a$ with $f_a\in A$. The \emph{Frobenius trace} is
$\Phi_u(f)=f_{(p-1,\dots,p-1)}$, and $\Phi^e_u$ is its $e$-th iterate. For two $p$-bases $u,u'$ of a field, $J(u/u')=\det(\partial u_i/\partial u'_j)$,
where the $\partial/\partial u'_j$ are the derivations attached to $u'$ ($\partial u'_i/\partial u'_j=\delta_{ij}$). For $X=\PP^n$, the variables of
$k[x,t]$ form a $p$-basis and $\Phi^e$ is the trace of Lemma~\ref{lem:frobenius-roots}.
\end{convention}

\begin{lemma}[Basic properties of $p$-bases {\cite[Lemma~3.3]{FHX}}]\label{lem:sg-pbasis-basics}
\lean{EnvelopeProperty.Cartier.IsPBasis.sum_traceIter_mul_mono, EnvelopeProperty.Cartier.IsPBasis.traceIter_mono_sub_mul_sum, EnvelopeProperty.Cartier.IsPBasis.traceIter_eq_of_sum_eq, EnvelopeProperty.Cartier.IsPBasis.traceIter_mono_pow_sub_one, EnvelopeProperty.Cartier.IsPBasis.traceIter_eq_traceIter_jac_pow_mul, EnvelopeProperty.Cartier.IsPBasis.traceIter_eq_traceIter_jac_pow_mul_of_isDomain, EnvelopeProperty.Cartier.IsPBasis.map_algebraMap_of_isLocalization, EnvelopeProperty.Cartier.IsPBasis.traceIter_mk', EnvelopeProperty.Cartier.IsPBasis.polynomial, EnvelopeProperty.Cartier.IsPBasis.laurentPolynomial, EnvelopeProperty.Cartier.IsPBasis.polynomial_laurentPolynomial}
\leanok
\uses{conv:sg-pbasis, lem:pn-divisorial-trace}
Let $u$ be a $p$-basis of $A$ and $Q=p^e$.
\begin{enumerate}
\item Every $f\in A$ can be written uniquely as $f=\sum_{a\in[0,Q)^N}f_{[a]}^Qu^a$; then $\Phi^e_u(f)=f_{[Q-1,\dots,Q-1]}$ and
$f_{[a]}=\Phi^e_u(u^{Q-1-a}f)$. The map $\Phi^e_u$ is additive, $\Phi^e_u(g^Qf)=g\Phi^e_u(f)$, $\Phi^e_u(u^{Q-1})=1$ and
$\Phi^{e'}_u\circ\Phi^e_u=\Phi^{e+e'}_u$.
\item If $A$ is a domain, $u$ is a $p$-basis of $\operatorname{Frac}A$, and the same formulas define $\Phi^e_u$ there. If $u'$ is another $p$-basis of
a field $F$, then $\Phi^e_{u'}(f)=\Phi^e_u(J(u/u')^{Q-1}f)$ for $f\in F$, and $J(u/u'')=J(u/u')J(u'/u'')$ for a third $p$-basis $u''$. If $u$ and
$u'$ are $p$-bases of a domain $A$, then $J(u/u')$ is a unit of $A$.
\item For $h\in A$, $u$ is a $p$-basis of $A_h$, and $\Phi^e_u(f/h^l)=\Phi^e_u(h^{l(Q-1)}f)/h^l$.
\item $(u,Y,t)$ is a $p$-basis of $A[Y,Y^{-1}][t]$.
\end{enumerate}
\end{lemma}

\begin{proof}
(1) Induction on $e$: if $f=\sum_bf_b^pu^b$ and $f_b=\sum_cf_{bc}^{Q/p}u^c$ with $c\in[0,Q/p)^N$, then $f=\sum_{b,c}f_{bc}^Qu^{pc+b}$, and
$(b,c)\mapsto pc+b$ is a bijection onto $[0,Q)^N$; uniqueness follows the same way. The top index $Q-1=p(Q/p-1)+(p-1)$ gives
$\Phi^e_u=\Phi^{e-1}_u\circ\Phi_u$. Semilinearity follows from uniqueness. For $f_{[a]}$: $u^{Q-1-a}f=\sum_bf_{[b]}^Qu^{Q-1-a+b}$, and writing
$Q-1-a+b=Q\gamma+\rho$ with $\rho\in[0,Q)^N$, we have $\rho=Q-1$ if and only if $b=a$ (as $|b_i-a_i|<Q$), in which case $\gamma=0$;
so $\Phi^e_u(u^{Q-1-a}f)=f_{[a]}$.
(2) The first claim holds because $f/g=fg^{p-1}/g^p$. The chain rule for the attached derivations gives $J(u/u'')=J(u/u')J(u'/u'')$.
For the transformation rule, let $C$ be the Cartier operator on top forms of $F$, which does not depend on the $p$-basis and satisfies
$C(g\,du)=\Phi_u(g)\,du$ (Lemma~\ref{lem:pn-divisorial-trace}). With $J=J(u/u')$ we have $du=J\,du'$, so
$\Phi_{u'}(f)\,du'=C(f\,du')=C(fJ^{-1}du)=\Phi_u(fJ^{-1})\,J\,du'$, and $J\Phi_u(fJ^{-1})=\Phi_u(J^{p-1}f)$ by semilinearity. If
$\Phi^e_{u'}(g)=\Phi^e_u(J^{Q-1}g)$ for all $g$, then
$\Phi^{e+1}_{u'}(f)=\Phi^e_u(J^{Q-1}\Phi_u(J^{p-1}f))=\Phi^e_u(\Phi_u(J^{pQ-p}J^{p-1}f))=\Phi^{e+1}_u(J^{pQ-1}f)$, which gives the formula by induction.
If $u,u'$ are $p$-bases of $A$, the derivations attached to either preserve $A$, so $J(u/u'),J(u'/u)\in A$, and their product is $J(u/u)=1$.
(3) $f/h^l=(h^{l(p-1)}f)/h^{lp}$ and the expansion of $h^{l(p-1)}f$ give existence; uniqueness follows by clearing denominators. Apply (1).
(4) $A[Y^{\pm1}][t]$ is free over its subring of $p$-th powers with basis the $u^aY^bt^c$, $a\in[0,p)^N$, $b,c\in[0,p)$: write
$Y^\beta=(Y^{\lfloor\beta/p\rfloor})^pY^{\beta-p\lfloor\beta/p\rfloor}$ and similarly for $t$.
\end{proof}

\begin{definition}[Frobenius roots {\cite[(3.3)]{FHX}}]\label{def:sg-root}
\lean{EnvelopeProperty.Cartier.IsPBasis.frobRoot}
\leanok
\uses{conv:sg-pbasis, lem:sg-pbasis-basics}
Let $u$ be a $p$-basis of $A$ and $\mathfrak a\subseteq A$ an ideal. The \emph{Frobenius root} is $\mathfrak a^{[1/p^e]}=\Phi^e_u(\mathfrak a)$, an ideal of $A$
(it is additive, and $g\Phi^e_u(f)=\Phi^e_u(g^{p^e}f)$). $\mathfrak a^{[q]}$ denotes the ideal generated by the $q$-th powers of elements of $\mathfrak a$.
\end{definition}

\begin{lemma}[Properties of Frobenius roots {\cite[(3.3)]{FHX}}]\label{lem:sg-root-basic}
\lean{EnvelopeProperty.Cartier.IsPBasis.frobRoot_le_iff, EnvelopeProperty.Cartier.IsPBasis.frobRoot_frobPow, EnvelopeProperty.Cartier.IsPBasis.frobRoot_frobRoot, EnvelopeProperty.Cartier.IsPBasis.frobRoot_mul_le, EnvelopeProperty.Cartier.IsPBasis.frobRoot_mono, EnvelopeProperty.Cartier.IsPBasis.frobRoot_eq_frobRoot, EnvelopeProperty.Cartier.IsPBasis.frobRoot_map}
\leanok
\uses{def:sg-root, lem:sg-pbasis-basics}
Let $u$ be a $p$-basis of $A$, $Q=p^e$, $Q'=p^{e'}$, and $\mathfrak a,\mathfrak b\subseteq A$ ideals.
\begin{enumerate}
\item $\mathfrak a^{[1/Q]}\subseteq\mathfrak b$ if and only if $\mathfrak a\subseteq\mathfrak b^{[Q]}$.
\item $(\mathfrak a^{[Q]})^{[1/Q]}=\mathfrak a$ and $(\mathfrak a^{[1/Q]})^{[1/Q']}=\mathfrak a^{[1/QQ']}$.
\item $(\mathfrak a\mathfrak b)^{[1/Q]}\subseteq\mathfrak a^{[1/Q]}\mathfrak b^{[1/Q]}$, and $\mathfrak a\mapsto\mathfrak a^{[1/Q]}$ is increasing.
\end{enumerate}
\end{lemma}

\begin{proof}
(1) If $\mathfrak a\subseteq\mathfrak b^{[Q]}$ and $f=\sum_jb_j^Qc_j\in\mathfrak a$ with $b_j\in\mathfrak b$, then $\Phi^e_u(f)=\sum_jb_j\Phi^e_u(c_j)\in\mathfrak b$. Conversely every
$f\in\mathfrak a$ is $\sum_af_{[a]}^Qu^a$ with $f_{[a]}=\Phi^e_u(u^{Q-1-a}f)\in\mathfrak a^{[1/Q]}$ (Lemma~\ref{lem:sg-pbasis-basics}(1)), so
$\mathfrak a\subseteq(\mathfrak a^{[1/Q]})^{[Q]}$, which lies in $\mathfrak b^{[Q]}$ if $\mathfrak a^{[1/Q]}\subseteq\mathfrak b$.
(2) $\Phi^e_u(\sum_jb_j^Qc_j)=\sum_jb_j\Phi^e_u(c_j)$ shows $(\mathfrak a^{[Q]})^{[1/Q]}\subseteq\mathfrak a$, and $b=\Phi^e_u(b^Qu^{Q-1})$ the converse. The second
identity is $\Phi^{e'}_u\circ\Phi^e_u=\Phi^{e+e'}_u$.
(3) By (1), $\mathfrak a\subseteq(\mathfrak a^{[1/Q]})^{[Q]}$ and $\mathfrak b\subseteq(\mathfrak b^{[1/Q]})^{[Q]}$, so $\mathfrak a\mathfrak b\subseteq(\mathfrak a^{[1/Q]}\mathfrak b^{[1/Q]})^{[Q]}$; apply (1).
Monotonicity is clear.
\end{proof}

\begin{lemma}[Test elements]\label{lem:sg-test-element}
\lean{EnvelopeProperty.Cartier.IsPBasis.exists_frobRoot_span_singleton_eq_top}
\leanok
\uses{lem:sg-root-basic}
Let $A$ be a Noetherian domain with a $p$-basis and $0\ne g\in A$. Then $(g)^{[1/p^e]}=A$ for some $e$.
\end{lemma}

\begin{proof}
By Lemma~\ref{lem:sg-root-basic}(2), $(g)^{[1/p^e]}=((g^p))^{[1/p^{e+1}]}\subseteq(g)^{[1/p^{e+1}]}$, since $(g^p)=(g)^{[p]}\subseteq(g)$; so the ideals
$(g)^{[1/p^e]}$ increase and stabilize at some $\mathfrak c$. If $\mathfrak c\ne A$, let $\mathfrak m\supseteq\mathfrak c$ be maximal. For every $e$,
$(g)\subseteq((g)^{[1/p^e]})^{[p^e]}\subseteq\mathfrak m^{[p^e]}\subseteq\mathfrak m^{p^e}$ (Lemma~\ref{lem:sg-root-basic}(1)), so the image of $g$ in the Noetherian
local domain $A_{\mathfrak m}$ lies in $\bigcap_e\mathfrak m^{p^e}A_{\mathfrak m}=0$ (Krull's intersection theorem), and $g=0$, a contradiction.
\end{proof}

\begin{definition}[Test ideals {\cite[(3.3)]{FHX}}]\label{def:sg-test-ideal}
\lean{EnvelopeProperty.Cartier.IsPBasis.testIdeal, EnvelopeProperty.Cartier.IsPBasis.testChain}
\leanok
\uses{def:sg-root, lem:test-ideals}
Let $A$ be a Noetherian domain with a $p$-basis $u$, $\mathfrak a\subseteq A$ an ideal and $c\in\Q_{\ge0}$. The test ideal is
$\tau(\mathfrak a^c)=\bigcup_{e\ge0}(\mathfrak a^{\lceil cp^e\rceil})^{[1/p^e]}$.
\end{definition}

\begin{lemma}[Stabilization {\cite[(3.3)]{FHX}}]\label{lem:sg-test-stable}
\lean{EnvelopeProperty.Cartier.IsPBasis.testChain_mono, EnvelopeProperty.Cartier.IsPBasis.exists_testIdeal_eq, EnvelopeProperty.Cartier.IsPBasis.testIdeal_mono, EnvelopeProperty.Cartier.IsPBasis.le_testIdeal_one}
\leanok
\uses{def:sg-test-ideal, lem:sg-root-basic}
The ideals $(\mathfrak a^{\lceil cp^e\rceil})^{[1/p^e]}$ increase with $e$; hence $\tau(\mathfrak a^c)=(\mathfrak a^{\lceil cp^e\rceil})^{[1/p^e]}$ for all $e\gg0$, and
$\tau(\mathfrak a^c)$ is increasing in $\mathfrak a$. Moreover $\mathfrak a\subseteq\tau(\mathfrak a^1)$.
\end{lemma}

\begin{proof}
Put $N_e=\lceil cp^e\rceil$, so $N_{e+1}\le pN_e$. By Lemma~\ref{lem:sg-root-basic}(2),(3),
$(\mathfrak a^{N_e})^{[1/p^e]}=((\mathfrak a^{N_e})^{[p]})^{[1/p^{e+1}]}\subseteq(\mathfrak a^{pN_e})^{[1/p^{e+1}]}\subseteq(\mathfrak a^{N_{e+1}})^{[1/p^{e+1}]}$, using
$(\mathfrak a^{N_e})^{[p]}\subseteq\mathfrak a^{pN_e}\subseteq\mathfrak a^{N_{e+1}}$. An increasing chain of ideals in a Noetherian ring stabilizes. Monotonicity is
Lemma~\ref{lem:sg-root-basic}(3), and $\mathfrak a=(\mathfrak a^{[p^e]})^{[1/p^e]}\subseteq(\mathfrak a^{p^e})^{[1/p^e]}\subseteq\tau(\mathfrak a^1)$.
\end{proof}

\begin{lemma}[Powers and subadditivity {\cite[(3.4)]{FHX}}]\label{lem:sg-test-subadditive}
\lean{EnvelopeProperty.Cartier.IsPBasis.frobRoot_pow_le_frobRoot_pow_add, EnvelopeProperty.Cartier.IsPBasis.testIdeal_pow, EnvelopeProperty.Cartier.IsPBasis.testIdeal_add_le, EnvelopeProperty.Cartier.IsPBasis.testIdeal_nsmul_le}
\leanok
\uses{lem:sg-test-element, lem:sg-test-stable, lem:sg-root-basic}
Let $A$ be a Noetherian domain with a $p$-basis and $0\ne\mathfrak a\subseteq A$.
\begin{enumerate}
\item For $0\ne g\in\mathfrak a$, $r\in\N$ and $k$ with $(g^r)^{[1/p^k]}=A$: $(\mathfrak a^N)^{[1/Q]}\subseteq(\mathfrak a^{Np^k+r})^{[1/p^kQ]}$ for all $N$ and $Q$.
\item $\tau((\mathfrak a^r)^c)=\tau(\mathfrak a^{rc})$ for $r\in\N$, $c\in\Q_{\ge0}$.
\item $\tau(\mathfrak a^{c+d})\subseteq\tau(\mathfrak a^c)\tau(\mathfrak a^d)$ for $c,d\in\Q_{\ge0}$; hence $\tau(\mathfrak a^{lc})\subseteq\tau(\mathfrak a^c)^l$ for $l\ge1$.
\end{enumerate}
\end{lemma}

\begin{proof}
(1) Let $\mathfrak c=g^r(\mathfrak a^N)^{[p^k]}\subseteq\mathfrak a^{r+Np^k}$. By semilinearity, $\mathfrak c^{[1/p^k]}=\mathfrak a^N\cdot(g^r)^{[1/p^k]}=\mathfrak a^N$, so by
Lemma~\ref{lem:sg-root-basic}(2),(3), $(\mathfrak a^N)^{[1/Q]}=\mathfrak c^{[1/p^kQ]}\subseteq(\mathfrak a^{Np^k+r})^{[1/p^kQ]}$; $k$ exists by
Lemma~\ref{lem:sg-test-element}.
(2) The terms of the chains satisfy $((\mathfrak a^r)^{\lceil cQ\rceil})^{[1/Q]}\subseteq(\mathfrak a^{\lceil rcQ\rceil})^{[1/Q]}$, since $r\lceil cQ\rceil\ge\lceil rcQ\rceil$,
giving $\subseteq$. Conversely, by (1) with this $r$, $(\mathfrak a^{\lceil rcQ\rceil})^{[1/Q]}\subseteq(\mathfrak a^{\lceil rcQ\rceil p^k+r})^{[1/p^kQ]}$, and
$\lceil rcQ\rceil p^k+r\ge r\lceil cp^kQ\rceil$, so this lies in $((\mathfrak a^r)^{\lceil cp^kQ\rceil})^{[1/p^kQ]}\subseteq\tau((\mathfrak a^r)^c)$.
(3) Choose $k$ as in (1) for $r=1$ and put $Q'=p^kQ$. Then $\lceil cQ'\rceil+\lceil dQ'\rceil\le\lceil(c+d)Q'\rceil+1\le\lceil(c+d)Q\rceil p^k+1$, so by (1) and
Lemma~\ref{lem:sg-root-basic}(3), $(\mathfrak a^{\lceil(c+d)Q\rceil})^{[1/Q]}\subseteq(\mathfrak a^{\lceil cQ'\rceil}\mathfrak a^{\lceil dQ'\rceil})^{[1/Q']}\subseteq
(\mathfrak a^{\lceil cQ'\rceil})^{[1/Q']}(\mathfrak a^{\lceil dQ'\rceil})^{[1/Q']}\subseteq\tau(\mathfrak a^c)\tau(\mathfrak a^d)$. Induction on $l$ gives the last claim.
\end{proof}

\begin{lemma}[Asymptotic test ideals {\cite[(3.4)]{FHX}}]\label{lem:sg-asymptotic}
\lean{EnvelopeProperty.Cartier.IsPBasis.asymptoticTestIdeal, EnvelopeProperty.Cartier.IsPBasis.testIdeal_le_testIdeal_of_dvd, EnvelopeProperty.Cartier.IsPBasis.exists_asymptoticTestIdeal_eq, EnvelopeProperty.Cartier.IsPBasis.testIdeal_le_asymptoticTestIdeal, EnvelopeProperty.Cartier.IsPBasis.le_asymptoticTestIdeal, EnvelopeProperty.Cartier.IsPBasis.asymptoticTestIdeal_mul_le, EnvelopeProperty.Cartier.IsPBasis.asymptoticTestIdeal_map}
\leanok
\uses{lem:sg-test-subadditive, lem:sg-test-stable, lem:test-ideals}
Let $A$ be a Noetherian domain with a $p$-basis and $(J_r)_{r\ge1}$ nonzero ideals with $J_rJ_{r'}\subseteq J_{r+r'}$. For $m\ge1$:
\begin{enumerate}
\item if $r\mid r'$, then $\tau(J_r^{m/r})\subseteq\tau(J_{r'}^{m/r'})$;
\item there is $r_0$ such that $B_m(A;J_\bullet)\coloneqq\tau(J_r^{m/r})$ is the same ideal for all $r$ divisible by $r_0$, and it contains
$\tau(J_r^{m/r})$ for every $r$;
\item $J_m\subseteq B_m$ and $B_{ml}\subseteq B_m^l$ for $l\ge1$.
\end{enumerate}
\end{lemma}

\begin{proof}
(1) With $s=r'/r$, $J_r^s\subseteq J_{r'}$, so $\tau(J_r^{m/r})=\tau((J_r^s)^{m/r'})\subseteq\tau(J_{r'}^{m/r'})$ by Lemma~\ref{lem:sg-test-subadditive}(2) and
Lemma~\ref{lem:sg-test-stable}. (2) The chain $\tau(J_{j!}^{m/j!})$ increases by (1) and stabilizes for $j\ge j_0$; put $r_0=j_0!$. For any $r$ and
$j\ge\max(j_0,r)$, $\tau(J_r^{m/r})\subseteq\tau(J_{j!}^{m/j!})$, and if $r_0\mid r$ the reverse inclusion follows from
$\tau(J_{j!}^{m/j!})=\tau(J_{r_0}^{m/r_0})\subseteq\tau(J_r^{m/r})$. (3) $J_m\subseteq\tau(J_m^1)=\tau(J_m^{m/m})\subseteq B_m$ by
Lemma~\ref{lem:sg-test-stable} and (2); and for $r$ divisible enough,
$B_{ml}=\tau(J_r^{ml/r})\subseteq\tau(J_r^{m/r})^l=B_m^l$ by Lemma~\ref{lem:sg-test-subadditive}(3).
\end{proof}

\subsection{\texorpdfstring{$p$}{p}-bases of smooth algebras}\label{subsec:sg-smooth}

\begin{lemma}[$p$-bases of function fields]\label{lem:sg-field-pbasis}
\lean{EnvelopeProperty.Cartier.finrank_kaehler_eq_trdeg, EnvelopeProperty.Cartier.finrank_kaehler_eq_toNat_trdeg, EnvelopeProperty.Cartier.span_D_eq_top_of_adjoin_eq_top, EnvelopeProperty.Cartier.span_D_eq_top_of_algebra_adjoin_eq_top, EnvelopeProperty.Cartier.isPBasis_of_span_D, EnvelopeProperty.Cartier.isPBasis_of_span_D_of_trdeg, EnvelopeProperty.Cartier.isPBasis_of_linearIndependent_D, EnvelopeProperty.Cartier.finrank_pthPowers_eq, EnvelopeProperty.Cartier.exists_isPBasis_of_trdeg, EnvelopeProperty.Cartier.exists_isPBasis_of_adjoin_eq_top_of_trdeg, EnvelopeProperty.Cartier.IsPBasis.finrank_kaehler}
\leanok
\uses{conv:sg-pbasis}
Let $F$ be a finitely generated field extension of $k$ of transcendence degree $d$.
\begin{enumerate}
\item $\dim_F\Omega_{F/k}=d$, and $\Omega_{F/k}$ is spanned by the differentials of any family of field generators of $F$ over $k$.
\item If $w_1,\dots,w_d\in F$ and $dw_1,\dots,dw_d$ is a basis of $\Omega_{F/k}$, then $w$ is a $p$-basis of $F$.
\item $[F:F^p]=p^d$.
\end{enumerate}
\end{lemma}

\begin{proof}
(1) Since $k$ is perfect, $F$ has a separating transcendence basis $y_1,\dots,y_d$, and $F/k(y)$ is finite separable, so
$\Omega_{F/k}=F\otimes\Omega_{k(y)/k}$ has basis $dy_1,\dots,dy_d$. The differential of a quotient is a combination of the differentials of numerator
and denominator, so field generators have spanning differentials.
(2) The $w_i$ are algebraically independent: otherwise take $0\ne P\in k[W_1,\dots,W_d]$ of minimal degree with $P(w)=0$; then
$\sum_i(\partial_iP)(w)\,dw_i=0$ forces $(\partial_iP)(w)=0$, hence $\partial_iP=0$ for all $i$ by minimality, so $P\in k[W^p]=k[W]^p$ ($k$ perfect),
$P=G^p$ with $G(w)=0$ and $\deg G<\deg P$, a contradiction. So $F$ is algebraic, hence finite, over $k(w)$, and $\Omega_{F/k(w)}=0$, so $F/k(w)$
is separable. The $w^a$, $a\in[0,p)^d$, form a basis of $k(w)$ over $k(w)^p=k(w^p)$. For a finite separable extension $L/K$ one has
$L=L^pK$ (the extension $L/L^pK$ is purely inseparable and separable), so a $p$-basis of $K$ spans $L$ over $L^p$; and $[L:L^p]=[K:K^p]$, since
Frobenius maps $L$ isomorphically onto $L^p$ and $K$ onto $K^p$, so $[L^p:K^p]=[L:K]$ and $[L:K][K:K^p]=[L:K^p]=[L:L^p][L^p:K^p]$. Hence the $p^d$
monomials $w^a$ form a basis of $F$ over $F^p$, which is the claim.
(3) follows from (1) and (2).
\end{proof}

\begin{lemma}[Local $p$-bases of smooth algebras {\cite[Lemma~3.3]{FHX}}]\label{lem:sg-kunz}
\lean{EnvelopeProperty.Cartier.exists_isPBasis_away_of_isSmoothAt, EnvelopeProperty.Cartier.exists_isPBasis_away_of_smooth, EnvelopeProperty.Cartier.exists_isPBasis_away_of_smooth_of_trdeg, EnvelopeProperty.Cartier.charP_of_isLocalization_away, EnvelopeProperty.Cartier.surjective_pSum_of_span_eq_maximalIdeal}
\leanok
\uses{lem:sg-field-pbasis, lem:sg-pbasis-basics}
Let $A$ be a smooth $k$-algebra of finite type which is a domain, $d=\operatorname{trdeg}_k\operatorname{Frac}A$, and $\mathfrak m\subseteq A$ a maximal ideal. There are
$u_1,\dots,u_d\in\mathfrak m$ and $h\in A\setminus\mathfrak m$ such that $u$ is a $p$-basis of $A_h$.
\end{lemma}

\begin{proof}
\emph{The cotangent space.} $A/\mathfrak m=k$ (Nullstellensatz). Since $A$ is smooth, $\Omega_{A/k}$ is projective, and its rank is
$\dim_K\Omega_{K/k}=d$ for $K=\operatorname{Frac}A$ (Lemma~\ref{lem:sg-field-pbasis}(1)). The conormal map $\mathfrak m/\mathfrak m^2\to\Omega_{A/k}\otimes_Ak$,
$f\mapsto df\otimes1$, is injective, because $A/\mathfrak m^2\to k$ has a $k$-algebra section, and surjective, because $da=d(a-\bar a)$ for the
constant $\bar a\in k$ with $a-\bar a\in\mathfrak m$. Hence $\dim_k\mathfrak m/\mathfrak m^2=d$. Choose $u_1,\dots,u_d\in\mathfrak m$ whose classes form a basis and put
$R=A_{\mathfrak m}$; by Nakayama, $\mathfrak mR=(u_1,\dots,u_d)$.

\emph{Spanning.} If $A=k[a_1,\dots,a_s]$, then $A=\sum_{\beta\in[0,p)^s}A^pa^\beta$ ($k=k^p$), and $R=\sum_\beta R^pa^\beta$ since $f/g=fg^{p-1}/g^p$; so
$R$ is a finite module over the local ring $R^p$, whose maximal ideal $\{f^p:f\in\mathfrak mR\}$ generates $(u_1^p,\dots,u_d^p)R$ in $R$ (as
$f=\sum r_iu_i$ gives $f^p=\sum r_i^pu_i^p$). Since $\mathfrak mR=(u)$ and $R/\mathfrak mR=k$, $R/\mathfrak m^jR$ is spanned over $k$ by the monomials $u^\beta$
with $|\beta|<j$; for $j=d(p-1)+1$ every monomial of degree $j$ has an exponent $\ge p$, so $\mathfrak m^jR\subseteq(u^p)$ and $R/(u^p)$ is spanned by
the $u^a$, $a\in[0,p)^d$. Nakayama over $R^p$ gives $R=\sum_aR^pu^a$.

\emph{Independence.} The $u^a$ span $K$ over $K^p=\operatorname{Frac}(R^p)$, and $[K:K^p]=p^d$ (Lemma~\ref{lem:sg-field-pbasis}(3)), so they are linearly
independent over $K^p$, hence over $R^p$; thus $u$ is a $p$-basis of $R$.

\emph{Spreading out.} Consider the $A^p$-linear map $\psi\colon(A^p)^{[0,p)^d}\to A$, $(c_a)\mapsto\sum_ac_au^a$. Localizing at the multiplicative set
$\{g^p:g\notin\mathfrak m\}$ turns $A^p$ into $R^p$ and $A$ into $R$, and $\psi$ into an isomorphism. Kernel and cokernel of $\psi$ are finitely
generated $A^p$-modules ($A^p\cong A$ is Noetherian, $A$ is finite over $A^p$), so a single $g^p$, $g\notin\mathfrak m$, kills both; then $\psi$ is an
isomorphism after inverting $g$, i.e.\ $u$ is a $p$-basis of $A_g$. Take $h=g$.
\end{proof}

\begin{corollary}[Covering by $p$-bases]\label{cor:sg-pbasis-cover}
\lean{EnvelopeProperty.Cartier.exists_isPBasis_cover_of_smooth, EnvelopeProperty.Cartier.exists_isPBasis_cover_of_smooth_of_trdeg}
\leanok
\uses{lem:sg-kunz}
Let $A$ be a smooth $k$-algebra of finite type which is a domain, $d=\operatorname{trdeg}_k\operatorname{Frac}A$. There are finitely many $h_1,\dots,h_s\in A$
generating the unit ideal and $p$-bases of $d$ elements of each $A_{h_\rho}$.
\end{corollary}

\begin{proof}
The elements $h$ provided by Lemma~\ref{lem:sg-kunz} for all maximal ideals generate an ideal contained in no maximal ideal, hence the unit ideal;
finitely many of them suffice.
\end{proof}

\subsection{Test ideals on the cone over a smooth projectively normal variety}\label{subsec:sg-cone}

In this subsection $X$ is projectively normal, has smooth charts, and $n\ge1$.

\begin{lemma}[Projective normality and the punctured cone]\label{lem:sg-proj-normal-sat}
\lean{EnvelopeProperty.ProjSubvariety.dvd_of_forall_dvd_mul_pow, EnvelopeProperty.ProjSubvariety.exists_eq_algebraMap_of_forall_mul_pow}
\leanok
\uses{lem:sg-cone-charts, lem:bj-dvr-minimal, conv:pg-setting}
$\bigcap_{i\in\mathcal J}S_{\bar x_i}=S$ in $\operatorname{Frac}S$, hence $\bigcap_{i\in\mathcal J}T_{\bar x_i}=T$.
\end{lemma}

\begin{proof}
Let $f$ lie in the intersection and fix $j\in\mathcal J$; write $f=g/h$ with $g\in S$ and $h=\bar x_j^N$. As $S$ is a Noetherian integrally closed
domain, by Lemma~\ref{lem:bj-dvr-minimal} it suffices that $\operatorname{ord}_P(g)\ge\operatorname{ord}_P(h)$ for every minimal prime $P$ of $hS$, at which $S_P$ is a discrete
valuation ring. Suppose $\mathfrak m_+\subseteq P$. Then $P=\mathfrak m_+$, which is maximal as $S/\mathfrak m_+=k$, and $S_{\mathfrak m_+}$ has dimension $1$. Since
$n\ge1$, the finitely generated domain $A_j$ is not a field (a field finitely generated as a $k$-algebra is algebraic over $k$), so it has
a nonzero maximal ideal $\mathfrak n$. Through $S\subseteq S_{\bar x_j}\cong A_j[Y^{\pm1}]$ (Lemma~\ref{lem:sg-cone-charts}), $\mathfrak p=S\cap\mathfrak nA_j[Y^{\pm1}]$ is a
prime ideal ($(A_j/\mathfrak n)[Y^{\pm1}]$ is a domain); it consists of the $\bar f=\sum_c\bar f_c$ with all $\bar f_c/\bar x_j^c\in\mathfrak n$, so it is homogeneous
with zero degree-$0$ part, hence $\mathfrak p\subseteq\mathfrak m_+$; it is nonzero (it contains $a\bar x_j^c$ for $0\ne a\in\mathfrak n$ of the form
$\bar g'/\bar x_j^c$) and does not contain $\bar x_j$. The chain $0\subsetneq\mathfrak p\subsetneq\mathfrak m_+$ contradicts $\dim S_{\mathfrak m_+}=1$. So some
$\bar x_i\notin P$, $i\in\mathcal J$, and $f\in S_{\bar x_i}\subseteq S_P$ gives $\operatorname{ord}_P(g)-\operatorname{ord}_P(h)=\operatorname{ord}_P(f)\ge0$. The statement for $T$ follows
coefficientwise in $t$.
\end{proof}

\begin{lemma}[A homogeneous $p$-basis]\label{lem:sg-homog-pbasis}
\lean{EnvelopeProperty.ProjSubvariety.coneZ, EnvelopeProperty.ProjSubvariety.exists_isPBasis_coneZ, EnvelopeProperty.ProjSubvariety.IsBihom, EnvelopeProperty.ProjSubvariety.IsBihom.traceIter_coneZ, EnvelopeProperty.ProjSubvariety.exists_eq_coneMk_of_isBihom, EnvelopeProperty.Cartier.IsPBasis.of_isUnit_jac, EnvelopeProperty.Cartier.IsPBasis.traceIter_laurentFamily_monomial}
\leanok
\uses{lem:sg-field-pbasis, lem:sg-cone-charts, lem:sg-pbasis-basics}
Fix $i_0\in\mathcal J$, put $Y=\bar x_{i_0}$ and $L=k(X)[Y,Y^{-1}][t]$, bigraded with $L_{(b,\lambda)}=k(X)Y^bt^\lambda$, so that
$T\subseteq T_{\bar x_{i_0}}\subseteq L\subseteq K\coloneqq\operatorname{Frac}T$ compatibly with the bigradings (Lemma~\ref{lem:sg-cone-charts}).
\begin{enumerate}
\item There are $i_1,\dots,i_n\in\mathcal J$ such that $w=(\bar x_{i_1}/\bar x_{i_0},\dots,\bar x_{i_n}/\bar x_{i_0})$ is a $p$-basis of $k(X)$. Then
$\zeta=(w,Y,t)$ is a $p$-basis of $L$ and of $K$, and $z=(\bar x_{i_0},\bar x_{i_1},\dots,\bar x_{i_n},t)$ is a $p$-basis of $K$ consisting of
bihomogeneous elements of $T$, with $\Phi^e_z(f)=\Phi^e_\zeta(Y^{-n(p^e-1)}f)$.
\item $\Phi^e_z(L)\subseteq L$, and for $f\in L_{(b,\lambda)}$ and $Q=p^e$: $\Phi^e_z(f)\in L_{(b',\lambda')}$ with
$b'=(b-(n+1)(Q-1))/Q$, $\lambda'=(\lambda-(Q-1))/Q$ if these are integers, and $\Phi^e_z(f)=0$ otherwise.
\end{enumerate}
\end{lemma}

\begin{proof}
(1) $k(X)$ is generated over $k$ by the $\bar x_j/\bar x_{i_0}$, so by Lemma~\ref{lem:sg-field-pbasis}(1),(2) some $n$ of them have differentials forming a
basis and are a $p$-basis; each $\bar x_{i_l}/\bar x_{i_0}$ is nonconstant, so $i_l\in\mathcal J$. By Lemma~\ref{lem:sg-pbasis-basics}(4) (applied to the
Laurent polynomial ring over $k(X)$), $\zeta$ is a $p$-basis of $L$, hence of $K$ (Lemma~\ref{lem:sg-pbasis-basics}(2)). The Jacobian matrix of
$z=(Y,w_1Y,\dots,w_nY,t)$ with respect to $\zeta$ is block triangular with determinant $\pm Y^n$, a unit of $K$; so $dz$ is a basis of $\Omega_{K/k}$
and $z$ is a $p$-basis of $K$ (Lemma~\ref{lem:sg-field-pbasis}(2)). By Lemma~\ref{lem:sg-pbasis-basics}(2),
$\Phi^e_z(f)=\Phi^e_\zeta(J(\zeta/z)^{Q-1}f)$ with $J(\zeta/z)=J(z/\zeta)^{-1}=\pm Y^{-n}$, and $(\pm1)^{Q-1}=1$ in characteristic $p$.
(2) For $f=cY^bt^\lambda$ with $c=\sum_{a\in[0,Q)^n}c_{[a]}^Qw^a\in k(X)$, writing $b=Q\lfloor b/Q\rfloor+\beta$ and $\lambda=Q\lfloor\lambda/Q\rfloor+\gamma$ gives
$f=\sum_a(c_{[a]}Y^{\lfloor b/Q\rfloor}t^{\lfloor\lambda/Q\rfloor})^Qw^aY^\beta t^\gamma$, so $\Phi^e_\zeta(f)=c_{[Q-1]}Y^{\lfloor b/Q\rfloor}t^{\lfloor\lambda/Q\rfloor}$ if
$\beta=\gamma=Q-1$ and $0$ otherwise. Combine with (1).
\end{proof}

\begin{definition}[Charts and Jacobians]\label{def:sg-charts}
\lean{EnvelopeProperty.ProjSubvariety.ChartDatum, EnvelopeProperty.ProjSubvariety.ChartDatum.isPBasis_v, EnvelopeProperty.ProjSubvariety.ChartDatum.jac, EnvelopeProperty.ProjSubvariety.exists_chartDatum_cover_all, EnvelopeProperty.ProjSubvariety.ConeTraceData}
\leanok
\uses{cor:sg-pbasis-cover, lem:sg-cone-charts, lem:sg-pbasis-basics, lem:sg-homog-pbasis}
Fix $i_0$ and $z$ as in Lemma~\ref{lem:sg-homog-pbasis}. For each $i\in\mathcal J$, Corollary~\ref{cor:sg-pbasis-cover} applied to the smooth domain
$A_i$ (of transcendence degree $n$) gives $h_{i\rho}\in A_i$, $\rho\in P_i$ finite, generating the unit ideal, and $p$-bases $u^{(i\rho)}$ of
$(A_i)_{h_{i\rho}}$. Write $h_{i\rho}=\bar{\tilde h}_{i\rho}/\bar x_i^{e_{i\rho}}$ with $\tilde h_{i\rho}\notin I$ a form of degree $e_{i\rho}$, and put
$s_{i\rho}=\bar x_i\bar{\tilde h}_{i\rho}\in S$ (homogeneous of positive degree) and $D_{i\rho}=T_{s_{i\rho}}=(A_i)_{h_{i\rho}}[\bar x_i^{\pm1}][t]$, a
bigraded subring of $L$ with the $p$-basis $v^{(i\rho)}=(u^{(i\rho)},\bar x_i,t)$ (Lemma~\ref{lem:sg-pbasis-basics}(3),(4)). Put
$J_{i\rho}=J(z/v^{(i\rho)})=\det(\partial z_a/\partial v^{(i\rho)}_b)$.
\end{definition}

\begin{lemma}[Properties of the charts]\label{lem:sg-jacobians}
\lean{EnvelopeProperty.ProjSubvariety.ChartDatum.jac_eq, EnvelopeProperty.ProjSubvariety.ChartDatum.isBihom_jac, EnvelopeProperty.ProjSubvariety.ChartDatum.jac_ne_zero, EnvelopeProperty.ProjSubvariety.ChartDatum.div_jac_mem, EnvelopeProperty.ProjSubvariety.ConeTraceData.traceIter_inv_jac_mul, EnvelopeProperty.Cartier.IsPBasis.jac_coneFamily}
\leanok
\uses{def:sg-charts, lem:sg-pbasis-basics, lem:sg-proj-normal-sat}
\begin{enumerate}
\item $J_{i\rho}$ is a nonzero bihomogeneous element of $D_{i\rho}$ of bidegree $(n,0)$.
\item $J_{i\rho}/J_{i'\rho'}$ is a unit of $T_{s_{i\rho}s_{i'\rho'}}$.
\item $\mathfrak m_+^{N_0}\subseteq(s_{i\rho})_{i,\rho}S$ for some $N_0$, and $\bigcap_{i,\rho}D_{i\rho}=T$.
\item $\Phi^e_z(J_{i\rho}^{-1}a)=J_{i\rho}^{-1}\Phi^e_{v^{(i\rho)}}(a)$ for $a\in D_{i\rho}$.
\end{enumerate}
\end{lemma}

\begin{proof}
Write $D=D_{i\rho}$, $v=(u,\bar x_i,t)=v^{(i\rho)}$. (1) The derivations attached to $v$ preserve $D$, and $z\subseteq T\subseteq D$, so $J\in D$. On
$cY_i^\beta t^\gamma$ with $c\in(A_i)_h$ and $Y_i=\bar x_i$, the expansion in $v$ is the expansion of $c$ in $u$ times the digits of $\beta$ and $\gamma$,
so $\partial/\partial u_b$ has bidegree $(0,0)$, $\partial/\partial Y_i$ bidegree $(-1,0)$ and $\partial/\partial t$ bidegree $(0,-1)$. Now
$z_l=(\bar x_{i_l}/\bar x_i)Y_i$ ($0\le l\le n$) has bidegree $(1,0)$ and $z_{n+1}=t$. In the Jacobian matrix the $t$-row and $t$-column are zero except
for $\partial t/\partial t=1$, the entries $\partial z_l/\partial u_b$ have bidegree $(1,0)$ and the $\partial z_l/\partial Y_i$ bidegree $(0,0)$; each term of the
determinant has bidegree $(n,0)$. $J\ne0$ because $J(z/v)J(v/z)=1$ in $K$ (Lemma~\ref{lem:sg-pbasis-basics}(2)).
(2) $J(z/v')=J(z/v)J(v/v')$, and $v,v'$ are $p$-bases of the domain $T_{ss'}$ (Lemma~\ref{lem:sg-pbasis-basics}(3)), so $J(v/v')$ is a unit there.
(3) From $1=\sum_\rho a_\rho h_{i\rho}$ in $A_i$, multiplying by a power of $\bar x_i$ gives $\bar x_i^{N}\in(s_{i\rho})_\rho S$, and $\bar x_i=0$ for
$i\notin\mathcal J$. Since $D_{i\rho}=(T_{\bar x_i})_{h_{i\rho}}$ and the $h_{i\rho}$ generate the unit ideal of $A_i$, $\bigcap_\rho D_{i\rho}=T_{\bar x_i}$; apply
Lemma~\ref{lem:sg-proj-normal-sat}.
(4) By Lemma~\ref{lem:sg-pbasis-basics}(2) in $K$, $\Phi^e_z(f)=\Phi^e_v(J(v/z)^{Q-1}f)=\Phi^e_v(J^{-(Q-1)}f)$, so
$\Phi^e_z(J^{-1}a)=\Phi^e_v(J^{-Q}a)=J^{-1}\Phi^e_v(a)$.
\end{proof}

\begin{definition}[Twisted forms]\label{def:sg-forms}
\lean{EnvelopeProperty.ProjSubvariety.ConeTraceData.iota, EnvelopeProperty.ProjSubvariety.ConeTraceData.mu, EnvelopeProperty.ProjSubvariety.ConeTraceData.M0, EnvelopeProperty.ProjSubvariety.ConeTraceData.J0, EnvelopeProperty.ProjSubvariety.ConeTraceData.mu_mul_iota, EnvelopeProperty.ProjSubvariety.ConeTraceData.isBihom_mu, EnvelopeProperty.ProjSubvariety.ConeTraceData.isBihom_iota, EnvelopeProperty.ProjSubvariety.ConeTraceData.span_M0_eq, EnvelopeProperty.ProjSubvariety.ConeTraceData.map_J0_eq}
\leanok
\uses{lem:sg-jacobians}
By Lemma~\ref{lem:sg-jacobians}(1),(2) choose $N_1$ such that $\iota_{i\rho}\coloneqq s_{i\rho}^{N_1}J_{i\rho}\in T$ and, for all $(i,\rho)$, $(i',\rho')$,
$s_{i\rho}^{N_1}J_{i'\rho'}/J_{i\rho}\in D_{i'\rho'}$ and $s_{i\rho}^{N_1}J_{i\rho}/J_{i'\rho'}\in D_{i'\rho'}$. Put $\mu_{i\rho}=s_{i\rho}^{N_1}J_{i\rho}^{-1}\in L$,
$M_0=\sum_{i,\rho}T\mu_{i\rho}\subseteq L$ and $\mathfrak J_0=\sum_{i,\rho}T\iota_{i\rho}\subseteq T$. The $\mu_{i\rho}$ and $\iota_{i\rho}$ are bihomogeneous; let
$\delta_M$ and $\delta_{\mathfrak J}$ be the largest $x$-degrees among them. Then $M_0D_{i\rho}=J_{i\rho}^{-1}D_{i\rho}$, $\mathfrak J_0D_{i\rho}=J_{i\rho}D_{i\rho}$, and
$\mu_{i\rho}\iota_{i\rho}=s_{i\rho}^{2N_1}$.
\end{definition}

The module $M_0$ plays the role of the regular top forms divided by $dz$, and $\mathfrak J_0$ of its inverse; the last three statements are
immediate from the choice of $N_1$, since $s_{i\rho}$ is a unit of $D_{i\rho}$.

\begin{lemma}[Saturation through the charts]\label{lem:sg-sat}
\lean{EnvelopeProperty.ProjSubvariety.ConeTraceData.mem_saturation_iff, EnvelopeProperty.ProjSubvariety.ConeTraceData.nonempty}
\leanok
\uses{lem:sg-jacobians, conv:sg-setting}
For an ideal $\mathfrak b\subseteq T$, $\mathfrak b^{\mathrm{sat}}=\{f\in T:f\in\mathfrak bD_{i\rho}\text{ for all }i,\rho\}$.
\end{lemma}

\begin{proof}
If $\mathfrak m_+^rf\subseteq\mathfrak b$, then $f\in\mathfrak bD_{i\rho}$ since $s_{i\rho}\in\mathfrak m_+$ is a unit of $D_{i\rho}$. Conversely, if $f\in\mathfrak bD_{i\rho}=\mathfrak bT_{s_{i\rho}}$ for
all $(i,\rho)$, then $s_{i\rho}^lf\in\mathfrak b$ for some $l$ and all $(i,\rho)$, and $\mathfrak m_+^{r}\subseteq(s_{i\rho}^l)_{i,\rho}$ for large $r$ by
Lemma~\ref{lem:sg-jacobians}(3).
\end{proof}

\begin{lemma}[The global trace on the charts]\label{lem:sg-theta-local}
\lean{EnvelopeProperty.ProjSubvariety.ConeTraceData.traceSubmodule, EnvelopeProperty.ProjSubvariety.ConeTraceData.thetaK, EnvelopeProperty.ProjSubvariety.ConeTraceData.theta, EnvelopeProperty.ProjSubvariety.ConeTraceData.map_theta, EnvelopeProperty.ProjSubvariety.ConeTraceData.exists_eq_algebraMap_of_mem_thetaK, EnvelopeProperty.ProjSubvariety.ConeTraceData.map_theta_chart}
\leanok
\uses{def:sg-forms, lem:sg-jacobians, lem:sg-homog-pbasis, def:sg-root, lem:sg-pbasis-basics}
For an ideal $\mathfrak a\subseteq T$ and $e,N\ge0$, $\Theta^e_N(\mathfrak a)\coloneqq\Phi^e_z(\mathfrak a^NM_0)\cdot\mathfrak J_0$ is an ideal of $T$, and
$\Theta^e_N(\mathfrak a)D_{i\rho}=(\mathfrak a^ND_{i\rho})^{[1/p^e]}$, the Frobenius root in $D_{i\rho}$ with respect to $v^{(i\rho)}$.
\end{lemma}

\begin{proof}
For a $T$-submodule $\mathfrak n\subseteq L$, $\Phi^e_z(\mathfrak n)$ is a $T$-submodule of $L$ by semilinearity (Lemma~\ref{lem:sg-homog-pbasis}(2) for
$\Phi^e_z(L)\subseteq L$), and $\Phi^e_z(\mathfrak n)T_s=\Phi^e_z(\mathfrak nT_s)$ for $s\in T\setminus0$, by Lemma~\ref{lem:sg-pbasis-basics}(3). With $D=D_{i\rho}$ and
$J=J_{i\rho}$: $\mathfrak a^NM_0D=\mathfrak a^NJ^{-1}D$, so by Lemma~\ref{lem:sg-jacobians}(4), $\Phi^e_z(\mathfrak a^NM_0)D=J^{-1}\Phi^e_v(\mathfrak a^ND)$, and multiplying by
$\mathfrak J_0D=JD$ gives $\Theta^e_N(\mathfrak a)D=\Phi^e_v(\mathfrak a^ND)=(\mathfrak a^ND)^{[1/p^e]}$. In particular $\Theta^e_N(\mathfrak a)\subseteq\bigcap_{i,\rho}D_{i\rho}=T$
(Lemma~\ref{lem:sg-jacobians}(3)).
\end{proof}

\begin{definition}[Asymptotic test ideals on the cone {\cite[(3.4)]{FHX}}]\label{def:sg-test-X}
\lean{EnvelopeProperty.ProjSubvariety.ConeTraceData.chartTestIdeal, EnvelopeProperty.ProjSubvariety.ConeTraceData.testSeq}
\leanok
\uses{lem:sg-asymptotic, def:sg-charts}
Let $(J_r)_{r\ge1}$ be nonzero ideals of $T$ with $J_rJ_{r'}\subseteq J_{r+r'}$. For $m\ge1$ let
$B_m=\{f\in T:f\in B_m(D_{i\rho};J_\bullet D_{i\rho})\text{ for all }i,\rho\}$, with the asymptotic test ideals of Lemma~\ref{lem:sg-asymptotic} in the
rings $D_{i\rho}$ with their $p$-bases $v^{(i\rho)}$.
\end{definition}

\begin{lemma}[Properties of $B_m$ {\cite[(3.4)]{FHX}}]\label{lem:sg-test-X}
\lean{EnvelopeProperty.ProjSubvariety.ConeTraceData.exists_theta_le_testSeq, EnvelopeProperty.ProjSubvariety.ConeTraceData.map_testSeq, EnvelopeProperty.ProjSubvariety.ConeTraceData.le_testSeq, EnvelopeProperty.ProjSubvariety.ConeTraceData.testSeq_mul_le_saturation, EnvelopeProperty.ProjSubvariety.ConeTraceData.exists_phiIdeal_testSeq_eq_theta, EnvelopeProperty.ProjSubvariety.ConeTraceData.phiIdeal_le_testSeq, EnvelopeProperty.ProjSubvariety.ConeTraceData.phiIdeal_testSeq_mul_le}
\leanok
\uses{def:sg-test-X, lem:sg-theta-local, lem:sg-sat, lem:sg-asymptotic, lem:sg-test-stable, lem:sg-phi-basic}
\begin{enumerate}
\item For every $m\ge1$, every $r$ divisible enough and every $e\gg0$ (depending on $m$ and $r$), with $N=\lceil mp^e/r\rceil$ and
$\Theta=\Theta^e_N(J_r)$: $\Theta\subseteq B_m\subseteq\Theta^{\mathrm{sat}}$ and $B_mD_{i\rho}=B_m(D_{i\rho};J_\bullet D_{i\rho})$ for all $i,\rho$.
\item $J_m\subseteq B_m$ and $B_{ml}\subseteq(B_m^l)^{\mathrm{sat}}$ for $m,l\ge1$.
\item $\phi_{B_m}=\phi_\Theta$, $\phi_{J_m}\le\phi_{B_m}$ and $\phi_{B_{ml}}\le l\,\phi_{B_m}$ on $X^\an$.
\end{enumerate}
\end{lemma}

\begin{proof}
(1) The $J_rD_{i\rho}$ are nonzero ideals of the Noetherian domain $D_{i\rho}$ forming a graded sequence. By Lemma~\ref{lem:sg-asymptotic}(2) and
Lemma~\ref{lem:sg-test-stable} for each of the finitely many charts, for $r$ divisible enough and $e\gg0$,
$B_m(D_{i\rho})=((J_rD_{i\rho})^{N})^{[1/p^e]}=\Theta D_{i\rho}$ (Lemma~\ref{lem:sg-theta-local}). Then $\Theta\subseteq B_m$ by definition, $B_m\subseteq\Theta^{\mathrm{sat}}$
by Lemma~\ref{lem:sg-sat}, and $\Theta D_{i\rho}\subseteq B_mD_{i\rho}\subseteq B_m(D_{i\rho})$.
(2) $J_m\subseteq B_m(D_{i\rho})$ by Lemma~\ref{lem:sg-asymptotic}(3). If $f\in B_{ml}$, then for all $(i,\rho)$,
$f\in B_{ml}(D_{i\rho})\subseteq B_m(D_{i\rho})^l=(B_mD_{i\rho})^l=B_m^lD_{i\rho}$ by Lemma~\ref{lem:sg-asymptotic}(3) and (1); Lemma~\ref{lem:sg-sat} gives
$f\in(B_m^l)^{\mathrm{sat}}$.
(3) Lemma~\ref{lem:sg-phi-basic}(1),(2) and (1), (2).
\end{proof}

\begin{proposition}[Uniform generation {\cite[Lemma~3.2]{FHX}}]\label{prop:sg-uniform-generation}
\lean{EnvelopeProperty.ProjSubvariety.ConeTraceData.c0, EnvelopeProperty.ProjSubvariety.ConeTraceData.theta_eq_span_genSet, EnvelopeProperty.ProjSubvariety.ConeTraceData.exists_span_le_testSeq_le_saturation, EnvelopeProperty.ProjSubvariety.ConeTraceData.isFS_inv_mul_phiIdeal_testSeq}
\leanok
\uses{lem:sg-test-X, lem:sg-homog-pbasis, def:sg-forms, lem:sg-phi-basic, lem:sg-pbasis-basics}
Let $c_0=M-n+1+\max(\delta_{\mathfrak J},0)$. Let $d\ge1$, let $(J_r)$ be as in Definition~\ref{def:sg-test-X} with each $J_r$ generated by bihomogeneous
elements of $x$-degree $rd$, and let $m\ge1$. There is an ideal $\mathfrak g\subseteq T$, generated by finitely many bihomogeneous elements of $x$-degree
$\le md+c_0$, with $\mathfrak g\subseteq B_m\subseteq\mathfrak g^{\mathrm{sat}}$. Consequently $m^{-1}\phi_{B_m}\in\mathcal H(X,d+c_0/m)$, with rational weights.
\end{proposition}

\begin{proof}
Take $r$, $e$ as in Lemma~\ref{lem:sg-test-X}(1) with moreover $Q=p^e\ge rd+\max(\delta_M,0)$, and $\mathfrak g=\Theta=\Theta^e_N(J_r)$, $N=\lceil mQ/r\rceil$.
The ideal $J_r^N$ is generated by products $g$ of $N$ generators of $J_r$, bihomogeneous of $x$-degree $Nrd$. Every monomial $x^\beta t^\lambda$ is
$(x^{\beta'}t^{\lambda'})^Qx^at^b$ with $a\in[0,Q)^{M+1}$, $b\in[0,Q)$, and $k=k^Q$; so $T=\sum_{a,b}T^Q\bar x^at^b$ and
$J_r^NM_0=\sum T^Q\,g\mu_{i\rho}\bar x^at^b$. By semilinearity $\Phi^e_z(J_r^NM_0)=\sum T\,\Phi^e_z(g\mu_{i\rho}\bar x^at^b)$, and $\Theta$ is generated by the
elements $\Phi^e_z(g\mu_{i\rho}\bar x^at^b)\,\iota_{i'\rho'}\in\Theta\subseteq T$. By Lemma~\ref{lem:sg-homog-pbasis}(2) they are bihomogeneous of $x$-degree at
most
\[
\frac{Nrd+\delta_M+(M+1)(Q-1)-(n+1)(Q-1)}Q+\delta_{\mathfrak J}\le md+\frac{rd+\delta_M}Q+(M-n)+\delta_{\mathfrak J}\le md+c_0,
\]
using $N\le mQ/r+1$, $M\ge n$ and $Q\ge rd+\max(\delta_M,0)$. Finally $\Theta\ne0$, because $0\ne J_m\subseteq B_m\subseteq\Theta^{\mathrm{sat}}$, so
Lemma~\ref{lem:sg-phi-basic}(3) and Lemma~\ref{lem:sg-test-X}(3) give $\phi_{B_m}=\phi_\Theta\in\mathcal H(X,md+c_0)$; scale by $m^{-1}$ (FS4).
\end{proof}

\begin{lemma}[Frobenius trace estimate {\cite[Lemma~3.3]{FHX}}]\label{lem:sg-trace-estimate}
\lean{EnvelopeProperty.Cartier.exists_logDiscrepancy}
\leanok
\uses{conv:sg-pbasis, lem:pn-divisorial-trace}
Let $R$ be a discrete valuation ring, essentially of finite type over $k$, with fraction field $F$, and $y$ a $p$-basis of $F$ consisting of elements
of $R$. There is $a\in\N$ with $\operatorname{ord}_R(\Phi^e_y(f))\ge\lfloor(\operatorname{ord}_R(f)+a)/p^e\rfloor-a$ for all $e\ge0$ and $f\in F$.
\end{lemma}

\begin{proof}
$R$ has a $p$-basis $u$ containing a uniformizer $\varpi$ (the case of a discrete valuation ring in Lemma~\ref{lem:pn-divisorial-trace}). Then $J=J(y/u)\in R$
is nonzero (Lemma~\ref{lem:sg-pbasis-basics}(2)); write $J=\epsilon\varpi^a$ with $\epsilon\in R^\times$. With $Q=p^e$, $\Phi^e_y(f)=\Phi^e_u(J^{-(Q-1)}f)=J^{-1}\Phi^e_u(Jf)$.
If $\operatorname{ord}_R(f)=o$, then $\operatorname{ord}_R(Jf)=o+a$, so $Jf=\varpi^{Q\ell}g$ with $\ell=\lfloor(o+a)/Q\rfloor$ and $g\in R$, and
$\Phi^e_u(Jf)=\varpi^\ell\Phi^e_u(g)\in\varpi^\ell R$ since $\Phi^e_u(R)\subseteq R$. Hence $\operatorname{ord}_R(\Phi^e_yf)\ge\ell-a$.
\end{proof}

\begin{proposition}[Valuative bound {\cite[Lemma~3.3, (3.8)]{FHX}}]\label{prop:sg-valuative-bound}
\lean{EnvelopeProperty.ProjSubvariety.ConeTraceData.ordGE_of_mem_thetaK, EnvelopeProperty.ProjSubvariety.ConeTraceData.exists_ord_testSeq_ge, EnvelopeProperty.ProjSubvariety.ConeTraceData.exists_phiIdeal_testSeq_le, EnvelopeProperty.ProjSubvariety.ConeTraceData.exists_phiIdeal_testSeq_coneObstacleIdeal_le}
\leanok
\uses{lem:sg-trace-estimate, lem:sg-test-X, def:sg-forms, lem:sg-homog-pbasis}
Let $R'$ be a discrete valuation ring with $T\subseteq R'\subseteq K$, essentially of finite type over $k$, such that $\bar x_j$ is a unit of $R'$ for
some $j\in\mathcal J$. There is $A\ge0$ such that for every graded sequence $(J_r)$ as in Definition~\ref{def:sg-test-X} and every $m\ge1$,
$\operatorname{ord}_{R'}(B_m)\ge\frac mr\operatorname{ord}_{R'}(J_r)-A$ for all sufficiently divisible $r$.
\end{proposition}

\begin{proof}
Let $a$ be given by Lemma~\ref{lem:sg-trace-estimate} for $R'$ and the $p$-basis $z\subseteq T\subseteq R'$ of $K$, and $b=\max_{i,\rho}\operatorname{ord}_{R'}(\iota_{i\rho})$; put
$A=1+a+b$. Since $\mu_{i\rho}\iota_{i\rho}=s_{i\rho}^{2N_1}\in T$, $\operatorname{ord}_{R'}(\mu_{i\rho})\ge-b$. Take $r$, $e$, $N$ and $\Theta$ as in
Lemma~\ref{lem:sg-test-X}(1). Every $f\in J_r^NM_0$ has $\operatorname{ord}_{R'}(f)\ge N\operatorname{ord}_{R'}(J_r)-b$, hence, with $Q=p^e$,
$\operatorname{ord}_{R'}(\Phi^e_zf)\ge(N\operatorname{ord}_{R'}(J_r)-b+a)/Q-1-a\ge\frac mr\operatorname{ord}_{R'}(J_r)-b-1-a$, as $N/Q\ge m/r$ and $\operatorname{ord}_{R'}(J_r)\ge0$. Elements of
$\mathfrak J_0\subseteq T$ have nonnegative order, so $\operatorname{ord}_{R'}(\Theta)\ge\frac mr\operatorname{ord}_{R'}(J_r)-A$. Finally $B_m\subseteq\Theta^{\mathrm{sat}}$ and
$B_m$ is finitely generated, so $\bar x_j^{l}B_m\subseteq\Theta$ for some $l$, and $\operatorname{ord}_{R'}(B_m)=\operatorname{ord}_{R'}(\bar x_j^lB_m)\ge\operatorname{ord}_{R'}(\Theta)$.
\end{proof}

\subsection{The envelope property}\label{subsec:sg-ep}

\begin{theorem}[Envelope property of smooth projective varieties in characteristic $p$ {\cite[Theorem~3.1]{FHX}}]\label{thm:sg-ep-charp}
\lean{EnvelopeProperty.ProjSubvariety.hasEnvelopeProperty_of_charP, EnvelopeProperty.ProjSubvariety.isPSH_envelope_nat_of_charP, EnvelopeProperty.ProjSubvariety.subsingleton_an_of_dim_eq_zero, EnvelopeProperty.ProjSubvariety.ConeTraceData.isPSH_subvarietyTestLimit, EnvelopeProperty.ProjSubvariety.ConeTraceData.subvarietyTestLimit_le_envelope_of_isDivisorial}
\leanok
\uses{thm:sg-reduction, lem:sg-pl-separate, prop:sg-envelope-Jm, lem:sg-Jm-basic, def:sg-test-X, lem:sg-test-X, prop:sg-uniform-generation, prop:sg-valuative-bound, def:sg-div, thm:sg-comparison, thm:pg-psh2, lem:sg-env-basic, rmk:bj-pn-ep-assembly}
Let $k$ be algebraically closed of characteristic $p>0$ and $X=V(I)\subseteq\PP^M$ projectively normal with smooth charts. Then every
$\theta\in\R$ has the envelope property on $X$.
\end{theorem}

\begin{proof}
If $n=0$, then $k(X)=k$ and every $A_i=k$, so each chart $U_{x_i}\cap X^\an\cong(\Spec k)^\an$ is the single point $\star$, $X^\an=\{\star\}$,
and the statement is clear ($\PSH(X,\theta)$ consists of the real constants if $\theta\ge0$ and is empty otherwise). Let $n\ge1$. By
Theorem~\ref{thm:sg-reduction} and Lemma~\ref{lem:sg-pl-separate} it suffices to show $\mathrm P_d(\phi)\in\PSH(X,d)$ for $d\in\mathbf{Z}_{>0}$ and
$\phi\in\operatorname{PL}(X)$ with $\phi\le0$. Fix the data of Definitions~\ref{def:sg-charts} and~\ref{def:sg-forms} and $c_0$ of
Proposition~\ref{prop:sg-uniform-generation}.

(a) Let $J_m=J_m(\phi)$ and $\mathrm Q=\mathrm P_d(\phi)=\sup_mm^{-1}\phi_{J_m}$ (Proposition~\ref{prop:sg-envelope-Jm}; $\phi$ is continuous). By
Lemma~\ref{lem:sg-Jm-basic}, $(J_m)$ is a graded sequence of nonzero ideals, $J_m$ generated by bihomogeneous elements of $x$-degree $md$. Let $B_m$ be
as in Definition~\ref{def:sg-test-X} and $u_m=m^{-1}\phi_{B_m}$, so $u_m\in\mathcal H(X,d+c_0/m)$ (Proposition~\ref{prop:sg-uniform-generation}) and $u_m\le0$.

(b) For $q,m\ge1$, Lemma~\ref{lem:sg-test-X}(3) and $J_q^m\subseteq J_{qm}$ give
$q^{-1}\phi_{J_q}\le(qm)^{-1}\phi_{J_{qm}}\le(qm)^{-1}\phi_{B_{qm}}\le m^{-1}\phi_{B_m}$; hence $\mathrm Q\le u_m$. Also $u_{(j+1)!}\le u_{j!}$. Let
$u=\lim_ju_{j!}$. Since $\mathrm Q\ge\inf\phi$ (Lemma~\ref{lem:sg-env-basic}), $u\ge\inf\phi>-\infty$, and PSH2 (Theorem~\ref{thm:pg-psh2}(3), classes
$d+c_0/j!\to d$) gives $u\in\PSH(X,d)$.

(c) Let $v\in X^{\divv}$ with $R'$, $c$ as in Definition~\ref{def:sg-div}, and $A$ as in Proposition~\ref{prop:sg-valuative-bound} (some $\bar x_j$ is a unit of $R'$).
For $m\ge1$ and $r$ divisible enough,
$\phi_{B_m}(v)=-c\operatorname{ord}_{R'}(B_m)\le-c\frac mr\operatorname{ord}_{R'}(J_r)+cA=\frac mr\phi_{J_r}(v)+cA\le m\mathrm Q(v)+cA$. Hence
$\mathrm Q(v)\le u(v)\le u_{j!}(v)\le\mathrm Q(v)+cA/j!$ for all $j$, and $u=\mathrm Q$ on $X^{\divv}$.

(d) Thus $u\le\phi$ on $X^{\divv}$; by Theorem~\ref{thm:sg-comparison} ($\phi$ is continuous) $u\le\phi$ on $X^\an$, so $u$ is a competitor and
$u\le\mathrm P_d(\phi)=\mathrm Q\le u$. Hence $\mathrm P_d(\phi)=u\in\PSH(X,d)$.
\end{proof}

\begin{remark}[Relation to the source]\label{rmk:sg-source}
\uses{thm:sg-ep-charp, rmk:bj-test-configurations, rmk:pg-subvariety-psh-comparison, prop:sg-uniform-generation, prop:sg-valuative-bound}
The proof of Theorem~\ref{thm:sg-ep-charp} is that of \cite[\S3]{FHX}, with the following differences.
\begin{itemize}
\item The ideals $J_m$ are defined from the obstacle (Definition~\ref{def:sg-Jm}) rather than from the base ideals of a test configuration
(Remark~\ref{rmk:bj-test-configurations}); \cite[Proposition~2.5(2)]{FHX} becomes Proposition~\ref{prop:sg-envelope-Jm}.
\item The test ideals live on the punctured affine cone over $X\times\A^1$, covered by the charts $D_{i\rho}$, rather than on
$S=X\times\A^1$; the cone is a $\mathbf G_m$-bundle over $S$, so these are the pullbacks of the test ideals of \cite{FHX}, but this
identification is not used. The canonical sheaf and the Grothendieck trace are replaced by the Frobenius trace $\Phi^e_z$ for the homogeneous
$p$-basis $z$ and the modules $M_0$, $\mathfrak J_0$, whose localizations are the regular top forms divided by $dz$ and their inverse.
\item The uniform twist of \cite[Lemma~3.2]{FHX} is $c_0\,\mathcal O(1)$ with $c_0$ of Proposition~\ref{prop:sg-uniform-generation}, obtained by counting
degrees, instead of $A=K_X+(n+1)H$, obtained from relative Castelnuovo--Mumford regularity and Serre vanishing; only the convergence of the
classes $d+c_0/m$ to $d$ matters.
\item The divisorial bound \cite[Lemma~3.3]{FHX} is used in the form of Proposition~\ref{prop:sg-valuative-bound}, with a constant $A=1+a+b$ that need
not be the log discrepancy; only its independence of $m$, $r$ and $J_\bullet$ matters.
\item $X$ is modelled as a projectively normal $V(I)$ with smooth charts and the classes as multiples of $\mathcal O(1)$
(Remark~\ref{rmk:pg-subvariety-psh-comparison}); projective normality is used only in Lemma~\ref{lem:sg-proj-normal-sat}.
\end{itemize}
For $X=\PP^n$ (with $I=0$), $z=(x_0,\dots,x_n,t)$ is a $p$-basis of $T=k[x,t]$ itself, its Frobenius roots are those of
Lemma~\ref{lem:frobenius-roots}, and Lemma~\ref{lem:sg-theta-local} holds with $M_0=\mathfrak J_0=T$, since each $J_{i\rho}$ is a unit of $D_{i\rho}=T_{x_i}$ (up to
sign a power of $x_i$); the argument is then that of Remark~\ref{rmk:bj-pn-ep-assembly}.
\end{remark}

\begin{theorem}[Characteristic-zero smooth envelopes from relative Nadel vanishing {\cite[Theorem~5.20]{BJ22}, \cite[Theorem~3.1]{FHX}}]\label{thm:import-sg-ep-char0}
\lean{EnvelopeProperty.smoothEnvelopeCharZero_of_relativeNadel}
\leanok
\uses{conv:sg-setting,def:sg-ep,def:c0n-relative-ample-nadel-input,thm:c0n-original-all-classes-envelope-from-nadel}
Assume the explicit relative ample ideal-power Nadel vanishing statement
of Definition~\ref{def:c0n-relative-ample-nadel-input}.
Let $k$ be algebraically closed of characteristic $0$ and let
$X=V(I)\subseteq\PP^M$ be projectively normal with smooth charts.
Then every $\theta\in\R_{\ge0}$ has the envelope property on $X$.
The conclusion is the original characteristic-zero smooth envelope
proposition, constructed in
Theorem~\ref{thm:c0n-original-smooth-envelope-from-nadel}.
\end{theorem}
\begin{proof}
Apply Theorem~\ref{thm:c0n-original-all-classes-envelope-from-nadel}
to the same original variety, smooth charts and class.
\end{proof}

\begin{theorem}[Envelope property of smooth projective varieties {\cite[Theorem~3.1]{FHX}}]\label{thm:sg-ep}
\lean{EnvelopeProperty.ProjSubvariety.hasEnvelopeProperty_of_smooth, EnvelopeProperty.ProjSubvariety.hasEnvelopeProperty_of_charZero, EnvelopeProperty.ProjSubvariety.hasEnvelopeProperty_of_neg}
\leanok
\uses{thm:sg-ep-charp, thm:import-sg-ep-char0, thm:c0-pn-ep}
Let $k$ be algebraically closed and $X=V(I)\subseteq\PP^M$ projectively normal with smooth charts. Then every $\theta\in\R$ has the envelope
property on $X$. In characteristic $0$ this follows from Theorem~\ref{thm:import-sg-ep-char0} under the explicit relative Nadel input, except for $X=\PP^n$ and $k$ countable,
where it is Theorem~\ref{thm:c0-pn-ep}. That input holds over every field of characteristic $0$
(Theorem~\ref{thm:kn-relative-ample-nadel-vanishing}); the statement without it is
Corollary~\ref{cor:kn-smooth-envelope-all-characteristics}.
\end{theorem}

\begin{proof}
In characteristic $p$ this is Theorem~\ref{thm:sg-ep-charp}. In characteristic $0$, for $\theta\ge0$ it is Theorem~\ref{thm:import-sg-ep-char0},
and for $\theta<0$ the class $\PSH(X,\theta)$ is empty (Proposition~\ref{prop:pg-maxprin}).
\end{proof}

\begin{corollary}[Smooth models of rank one {\cite[Theorem~3.1]{FHX}}]\label{cor:sg-smooth-model}
\lean{EnvelopeProperty.MultiModel.hasEnvelopeProperty_of_isSmooth, EnvelopeProperty.MultiModel.exists_degOnePres_veronese, EnvelopeProperty.MultiModel.DegOnePres.hasSmoothCharts_of_isSmooth}
\leanok
\uses{thm:sg-ep, def:gl-smooth, def:gl-psh, rmk:gl-model-examples, lem:gl-model-change, lem:gl-veronese, rmk:gl-bridges, lem:wu-veronese, lem:gl-charts}
Let $\mathcal Y=(K,S,z)$ be a smooth model of rank $1$ (Definition~\ref{def:gl-smooth}) and $t\ge0$. Then $\mathrm{EP}_{\mathcal Y}(t)$ holds. This is
Theorem~\ref{thm:gl-smooth-ep}; in characteristic $0$ it rests on Theorem~\ref{thm:import-sg-ep-char0}.
\end{corollary}

\begin{proof}
(a) $S_0\subseteq K$ is a domain integral over $k[z]_0=k$, hence a field algebraic over $k$, so $S_0=k$. Let $S^\nu$ be the integral closure of $S$ in
$K[T]$; by Remark~\ref{rmk:gl-model-examples}(2), $\mathcal Y^\nu=(K,S^\nu,z)$ is a normal model, and $S^\nu_0=k$ as before. The chart algebras of
$\mathcal Y^\nu$ contain those of $\mathcal Y$, are integral over them and have the same fraction field $K$; as the $A_j$ are integrally closed, they
coincide. By Lemma~\ref{lem:gl-model-change} restriction identifies $\mathcal Y^{\nu,\an}$ with $\mathcal Y^\an$ and the psh classes, so
$\mathrm{EP}_{\mathcal Y}(t)\iff\mathrm{EP}_{\mathcal Y^\nu}(t)$.

(b) By Lemma~\ref{lem:wu-veronese} there is $e\ge1$ such that $S^\nu_{(e)}=\bigoplus_nS^\nu_{ne}$ is generated by $S^\nu_e$ and integrally closed. Let
$w_0,\dots,w_N$ be a basis of $S^\nu_e$ and $I$ the kernel of $k[x_0,\dots,x_N]\to S^\nu_{(e)}$, $x_i\mapsto w_i$; then $X=V(I)\subseteq\PP^N$ is projectively
normal. By Lemma~\ref{lem:gl-veronese} and Remark~\ref{rmk:gl-bridges}, $\mathrm{EP}_{\mathcal Y^\nu}(et')\iff\mathrm{EP}_X(t')$ for $t'\ge0$.

(c) $X$ has smooth charts. For $w_i\ne0$ the chart algebra of $X$ at $x_i$ is $S^\nu_{(w_i)}$ (Lemma~\ref{lem:wu-veronese}). Since $w_i$ is integral over
$k[z]$, it satisfies $w_i^n+\sum_{l=1}^nc_lw_i^{n-l}=0$ with $c_l$ homogeneous of degree $le$ in $k[z]$; dividing by $w_i^n$ writes $1$ as a combination of
products of the elements $z^\beta/w_i$, $\beta$ a monomial of degree $e$, so these generate the unit ideal of $S^\nu_{(w_i)}$. The localization
of $S^\nu_{(w_i)}$ at $z^\beta/w_i$ is $S^\nu_{(w_iz^\beta)}$, a localization of $S^\nu_{(z^\beta)}$, which is the localization of the smooth chart algebra $A_j$
at $z^\beta/z_j^e$ for any $j$ with $\beta_j>0$. Smoothness is local for a finite cover by principal open sets, so $S^\nu_{(w_i)}$ is smooth.

(d) Theorem~\ref{thm:sg-ep} gives $\mathrm{EP}_X(t/e)$, hence $\mathrm{EP}_{\mathcal Y}(t)$ by (a) and (b).
\end{proof}

\begin{remark}[Characteristic zero]\label{rmk:sg-char0}
\uses{thm:import-sg-ep-char0, thm:kn-smooth-envelope-char-zero, thm:c0-pn-ep, rmk:c0-source, prop:c0-l2-extension, lem:c0-hormander-global, lem:c0-polynomial-growth}
Theorem~\ref{thm:import-sg-ep-char0} is the special case $\theta\,c_1(\mathcal O_X(1))$ of \cite[Theorem~5.20]{BJ22}
(Remark~\ref{rmk:pg-subvariety-psh-comparison}), which \cite[\S3]{FHX} cites in characteristic $0$. The proof of Theorem~\ref{thm:c0-pn-ep} for
$\PP^n$ uses that the cone over $\PP^n\times\A^1$ is the affine space $\mathbf C^{n+2}$: the Ohsawa--Takegoshi theorem at a point and
H\"ormander's estimate are applied on balls and on $\mathbf C^{n+2}$ with polynomial weights (Lemma~\ref{lem:c0-hormander-global},
Proposition~\ref{prop:c0-l2-extension}), and entire functions of polynomial growth are polynomials (Lemma~\ref{lem:c0-polynomial-growth}).
For general $X$ the cone $C(X)\times\mathbf C$ is singular along the vertex and is not a domain in a complex vector space. The corresponding
step needs an $L^2$ existence theorem for $\bar\partial$ on the complex manifold $X(\mathbf C)\times\PP^1(\mathbf C)$ with values in a line bundle
$\mathcal O(e,b)$ with positively curved, singular Hermitian metrics (H\"ormander's theorem on compact K\"ahler manifolds), or on the complete
K\"ahler manifold $(C(X)\setminus\{0\})\times\mathbf C$, together with the algebraicity of square integrable holomorphic sections. A finite linear
projection $C(X)\to\mathbf C^{n+2}$ does not reduce this to $\mathbf C^{n+2}$ directly: an estimate for a scalar weight on $\mathbf C^{n+2}$ controls the
$L^2$ norm on $C(X)$ only if the weight is dominated by the given weight on every sheet of the projection, and the fibrewise minimum of a
plurisubharmonic function is not plurisubharmonic in general. Finite descent \cite[\S4]{FHX} transports the envelope property from the source of
a finite morphism to its target, so a finite morphism $X\to\PP^n$ yields the envelope property of $\PP^n$ from that of $X$, and not
conversely.
\end{remark}

\begin{remark}[Nef classes]\label{rmk:sg-nef}
\uses{cor:sg-smooth-model, thm:gl-smooth-ep, cor:gl-smooth-ep-nef, lem:gl-ep-basic, def:gl-cones, rmk:gl-model-comparison, lem:sg-envppty}
\cite[Theorem~3.1]{FHX} concerns every nef class of a smooth projective variety; its reduction step chooses rational ample classes $\theta_j\to\theta$
with $\theta_j-\theta$ ample and applies \cite[Lemma~5.9]{BJ22}. In the model of \S\ref{sec:global} the nef classes are those of the model-nef cone
$N(\mathcal Y)$ (Remark~\ref{rmk:gl-model-comparison}). Corollary~\ref{cor:sg-smooth-model} is Theorem~\ref{thm:gl-smooth-ep}, and
Corollary~\ref{cor:gl-smooth-ep-nef} extends it to the classes $\theta\in\R^r_{\ge0}$ of a smooth model of rank $r$. For $\theta\in N(\mathcal Y)$ with a
negative coordinate, Lemma~\ref{lem:gl-ep-basic}(1) reduces the statement to the rational classes $a+\delta\mathbf 1$ with $a\in C(\mathcal Y)$ and
$\delta\in\Q_{>0}$ (Definition~\ref{def:gl-cones}), and these require a smooth model $\mathcal Y'$ of rank $1$ with the same analytification and
$\PSH(\mathcal Y',t)=\PSH(\mathcal Y,ts(a+\delta\mathbf 1))$ for some $s\in\mathbf Z_{>0}$ (a projective embedding by a multiple of the ample class
$a+\delta\mathbf1$), after which Corollary~\ref{cor:sg-smooth-model} applies. In a Fubini--Study system whose classes are compared coordinatewise,
the passage from $\Q^r_{>0}$ to $\R^r_{\ge0}$ is Lemma~\ref{lem:sg-envppty}.
\end{remark}
